% ============================================================
%  olowiane_dzieci_analiza.tex  —  Main file
%  Ołowiane dzieci. Zapomniana epidemia — Polish Study Companion
%  Compile with pdfLaTeX (or XeLaTeX — see font note below)
% ============================================================
%
% ── PROJECT CONTEXT (for session handoff) ───────────────────
%
%  This document is a phrase-by-phrase Polish language analysis
%  of the book "Ołowiane dzieci. Zapomniana epidemia" by Michał
%  Jędryk (a Polish-language book about a lead poisoning epidemic,
%  not available in English translation).
%
%  PURPOSE: Polish language learning through close reading.
%  The document is for private, personal, educational use only
%  and is not for distribution or commercial use.
%
% ── LEARNER PROFILE ─────────────────────────────────────────
%
%  - Native English speaker, age 61, based in Montreal, Canada,
%    with a condo in Kraków, Poland.
%  - Has ADHD and dyslexia/dysorthographia. LaTeX is an
%    accessibility aid: wheat background, clean typesetting,
%    structured navigation, and non-editable PDF output all
%    reduce visual fatigue. Avoid visually cluttered formatting.
%  - Documented working memory deficits (associated with ADHD):
%    rote memorisation is ineffective. Learning happens through
%    repeated contextual exposure — absorption, not drilling.
%  - Compensating strength (noted by neuropsychologist): strong
%    pattern recognition and rule derivation. Grammatical analysis
%    is not an academic exercise but the actual learning mechanism
%    — seeing the pattern once is worth more than memorising
%    paradigm tables. Frame analyses accordingly.
%  - Strong auditory learner — achieved near-native French
%    entirely by ear after immersion. Learns by osmosis rather
%    than rote memorisation.
%  - Latin background (high school): already understands noun
%    declension, verb conjugation, and grammatical cases
%    conceptually. Does not need these explained from scratch —
%    only the Polish instantiation of familiar concepts.
%  - Also has passive background in Russian and basic
%    conversational Polish from a Polish wife and in-laws.
%  - Goal: break through a conversational plateau; prepare for
%    regular extended stays in Kraków.
%
% ── ANALYSIS APPROACH ───────────────────────────────────────
%
%  - The user supplies screenshots of the book (DRM prevents
%    copying) one sentence or short passage at a time.
%  - The AI provides: full translation, word-by-word breakdown,
%    grammatical analysis, and literary/contextual commentary.
%  - Depth of unpacking is calibrated to interest — linger where
%    the language or content is rich, move faster elsewhere.
%  - Tangential grammar questions are welcomed and encouraged.
%  - The learner does NOT stress about memorisation; absorption
%    happens naturally through engagement with the content.
%  - LaTeX is updated periodically (not after every sentence) at
%    natural section breaks or roughly every 10–15 sentences.
%
% ── DOCUMENT STRUCTURE ──────────────────────────────────────
%
%  olowiane_dzieci_analiza.tex  — this file (main, compile this)
%  grammar.tex                  — Polish grammar foundations
%  ch00_karta_tytulowa.tex      — front matter passages
%  ch01_jeszcze_jeden_pogrzeb.tex — Chapter 1 (in progress)
%  ch02_... etc.                — subsequent chapters (stub)
%
%  All chapter files are \input{}'d below. Chapters not yet
%  started are listed but commented out.
%
% ── CUSTOM COMMANDS ─────────────────────────────────────────
%
%  \poltext{...}      — quoted Polish source text (blue italic)
%  \poltrans{...}     — English translation line
%  \polword{word}{gloss} — inline Polish term with translation
%  \grambox{title}{content} — shaded grammar note box
%
% ── USE OF AI ────────────────────────────────────────────────
%
%  All phrase-by-phrase analysis in this document was generated
%  with the assistance of Anthropic's Claude Sonnet 4.5 model
%  (model string: claude-sonnet-4-5), accessed via the Claude.ai
%  web interface. The analysis was produced interactively: the
%  learner supplied screenshots of the source text one passage
%  at a time, and the model provided translation, grammatical
%  breakdown, and literary commentary in a conversational
%  session. The LaTeX document was also structured, drafted, and
%  updated by the model in response to the learner's direction.
%
%  This workflow is described in full in the "Use of AI" preface
%  section of the compiled document, so that others may adapt
%  the method for their own language learning purposes.
%
%  The project is shared openly on GitHub. Neither the analysis
%  nor the document constitute a reproduction of the source text
%  — quoted passages are brief and used solely for educational
%  commentary consistent with fair use principles.
%
% ── SESSION HANDOFF INSTRUCTIONS ────────────────────────────
%
%  If you are a new AI session picking up this project:
%  1. Read this comment block in full.
%  2. Read the header comment block of the active chapter file
%     (currently ch01_jeszcze_jeden_pogrzeb.tex) — it contains
%     explicit style rules and the current reading position.
%  3. Read grammar.tex to understand the grammatical framework.
%  4. The user will supply passages as screenshots, one at a time.
%     Continue the phrase-by-phrase analysis in the same style.
%  5. Update the chapter .tex file after EACH passage is confirmed
%     by the user — do not batch updates. This is the agreed
%     protocol to avoid losing work between sessions.
%  6. Maintain the tone of existing analyses: informative,
%     literary where appropriate, and enjoyable to read in their
%     own right — not merely functional glosses.
%
%  CRITICAL STYLE RULES (details in chapter file header):
%  - ALL passage headings use \section{} — never \subsection.
%  - Polish source text always in \poltext{} — never \begin{quote}.
%  - English translation always in \poltrans{} — no label prefix.
%  - Analysis in free prose, \textbf{term} --- \textit{gloss} ---
%  - Read the chapter file header before writing any new content.
%
% ============================================================

\documentclass[12pt, a4paper, openany]{book}

% ── Encoding & Language ─────────────────────────────────────
\usepackage[T1]{fontenc}   % T1 required for Polish characters (ą ę ó ś ź ż ć ń ł)
\usepackage[utf8]{inputenc}
\usepackage[polish, english]{babel}

% ── Page geometry ───────────────────────────────────────────
\usepackage[
  a4paper,
  top=2.5cm, bottom=2.5cm,
  left=3cm,  right=3cm,
  marginparwidth=0cm,
  headheight=14pt
]{geometry}

% ── Colour ──────────────────────────────────────────────────
\usepackage{xcolor}

\definecolor{wheat}{RGB}{245, 222, 179}
\definecolor{polishblue}{RGB}{25, 80, 140}
\definecolor{analysisgrey}{RGB}{60, 60, 60}
\definecolor{keywordred}{RGB}{160, 30, 30}
\definecolor{rulecolour}{RGB}{180, 150, 100}
\definecolor{toccolour}{RGB}{25, 80, 140}

\pagecolor{wheat}
% body text is black on wheat — no override needed here

% ── Fonts ───────────────────────────────────────────────────
% If compiling with XeLaTeX/LuaLaTeX, use fontspec for best results.
% If using pdfLaTeX, the following provides reasonable output.
\usepackage{palatino}   % body font — good x-height, readable
\usepackage{microtype}  % improved kerning and spacing

% ── Spacing ─────────────────────────────────────────────────
\usepackage{setspace}
\setstretch{1.4}        % generous line spacing for readability
\setlength{\parskip}{0.6em}
\setlength{\parindent}{0pt}

% ── Headers & Footers ───────────────────────────────────────
\usepackage{fancyhdr}
\pagestyle{fancy}
\fancyhf{}
\fancyhead[LE,RO]{\small\color{polishblue}\itshape\leftmark}
\fancyfoot[C]{\small\color{analysisgrey}\thepage}
\renewcommand{\headrulewidth}{0.3pt}
\renewcommand{\headrule}{\color{rulecolour}\hrule width\headwidth height\headrulewidth}

% ── Hyperlinks & PDF metadata ────────────────────────────────
\usepackage[
  colorlinks=true,
  linkcolor=toccolour,
  urlcolor=toccolour,
  citecolor=toccolour,
  pdftitle={Ołowiane dzieci — Polish Study Companion},
  pdfauthor={Personal Study Document},
  bookmarks=true,
  bookmarksnumbered=true
]{hyperref}
\usepackage{xurl}          % allows long URLs to break at slashes

% ── Section title styling ────────────────────────────────────
\usepackage{titlesec}

\titleformat{\chapter}[display]
  {\normalfont\Large\bfseries\color{polishblue}}
  {\chaptertitlename\ \thechapter}{10pt}{\Huge}
\titlespacing{\chapter}{0pt}{0pt}{20pt}

\titleformat{\section}
  {\normalfont\large\bfseries\color{polishblue}}
  {}{0em}{}[\color{rulecolour}\titlerule]
\titlespacing{\section}{0pt}{1.5em}{0.8em}

\titleformat{\subsection}
  {\normalfont\normalsize\bfseries\color{keywordred}}
  {}{0em}{}
\titlespacing{\subsection}{0pt}{1em}{0.4em}

% ── TOC styling ──────────────────────────────────────────────
\usepackage{tocloft}
\renewcommand{\cftchapfont}{\bfseries\color{polishblue}}
\renewcommand{\cftchappagefont}{\bfseries\color{polishblue}}
\setlength{\cftbeforechapskip}{0.5em}
% Only show chapters in TOC (not sections/subsections)
\setcounter{tocdepth}{0}

% ── Custom environments and commands ─────────────────────────

% \poltext{} — display block for quoted Polish source text
\newcommand{\poltext}[1]{%
  \vspace{0.5em}
  {\color{keywordred}\large\itshape #1}%
  \vspace{0.3em}\par
}

% \poltrans{} — English translation line
\newcommand{\poltrans}[1]{%
  {\color{polishblue}\small\upshape #1}\par
  \vspace{0.5em}
  \color{rulecolour}\noindent\rule{\linewidth}{0.3pt}\color{black}
  \vspace{0.3em}
}

% \polword{word}{gloss} — inline Polish term with translation
\newcommand{\polword}[2]{%
  \textbf{\color{polishblue}#1}~{\color{keywordred}(#2)}%
}

% \grambox{title}{content} — shaded grammar note
\usepackage{mdframed}
\mdfdefinestyle{gramstyle}{%
  backgroundcolor=white!60!wheat,
  linecolor=rulecolour,
  linewidth=0.6pt,
  innertopmargin=6pt,
  innerbottommargin=6pt,
  innerleftmargin=10pt,
  innerrightmargin=10pt,
  skipabove=0.8em,
  skipbelow=0.8em
}
\newcommand{\grambox}[2]{%
  \begin{mdframed}[style=gramstyle]
    \textbf{\color{keywordred}#1}\quad #2
  \end{mdframed}
}

% ── Utility ──────────────────────────────────────────────────
\usepackage{enumitem}
\setlist[itemize]{leftmargin=1.5em, itemsep=0.2em}
\setlist[description]{leftmargin=1.5em, labelwidth=4em, itemsep=0.3em}

% ── Preface version selector ─────────────────────────────────
%  Change the value below to switch preface versions.
%  'short' — brief half-page AI disclosure (default)
%  'full'  — complete methodological preface (preface_ai.tex)
%  The \ifthenelse block in the body reads this and inputs the
%  correct file automatically — no other change needed.
\usepackage{ifthen}
\newcommand{\prefaceversion}{short}  % ← change to: short | full

% ── Fix default body colour after xcolor pagecolor ───────────
\color{black}
\makeatletter
\let\old@document\document
\renewcommand{\document}{\old@document\color{black}}
\makeatother

% ============================================================
\begin{document}
\color{black}  % ensure body text is black on wheat

% ── Title page ───────────────────────────────────────────────
\begin{titlepage}
  \centering
  \vspace*{3cm}
  {\Huge\bfseries\color{polishblue} Ołowiane dzieci}\par
  \vspace{0.5em}
  {\large\itshape Zapomniana epidemia}\par
  \vspace{0.3em}
  {\normalsize Michał Jędryk}\par
  \vspace{2.5cm}
  {\LARGE\color{keywordred} Polish Language Study Companion}\par
  \vspace{1em}
  {\normalsize A phrase-by-phrase reading with grammatical analysis}\par
  \vfill
  {\small\color{analysisgrey} Openly shared for educational purposes — see preface for licence and copyright notes}
\end{titlepage}

% ── Table of Contents ────────────────────────────────────────
\tableofcontents

% ── Preface: Use of AI ───────────────────────────────────────
%  Controlled by \prefaceversion in the preamble — edit that
%  value only; do not change anything here.
\ifthenelse{\equal{\prefaceversion}{full}}
  {% ============================================================
%  preface_ai.tex  —  Preface: Use of AI
%  Referenced by olowiane_dzieci_analiza.tex via % ============================================================
%  preface_ai.tex  —  Preface: Use of AI
%  Referenced by olowiane_dzieci_analiza.tex via % ============================================================
%  preface_ai.tex  —  Preface: Use of AI
%  Referenced by olowiane_dzieci_analiza.tex via \input{preface_ai}
% ============================================================

\chapter*{A Note on Method: Use of Artificial Intelligence}
\addcontentsline{toc}{chapter}{A Note on Method: Use of AI}
\markboth{A Note on Method}{}

This document was produced with the assistance of a large language
model (LLM) and is shared openly in the hope that others may adapt
the method for their own language learning.

\section*{The Learning Context}

The author is a native English speaker with a Polish wife, Polish
in-laws, and a recently acquired apartment in Kraków. After more than
twenty years of incidental exposure to Polish, conversational ability
had plateaued at a level sufficient for simple transactions but
inadequate for real two-way communication. The goal of this project
is to break through that plateau by engaging deeply with a Polish-language
book that is genuinely interesting, not available in English translation,
and whose subject matter --- a lead poisoning epidemic in industrial
Silesia --- connects directly to the author's professional background
in biomedical engineering.

The author has ADHD and dyslexia/dysorthographia, and shares this
openly in the hope of reducing the stigma that still surrounds these
conditions --- particularly in academic and language-learning contexts,
where they are often misread as indicators of low ability or
insufficient effort. \LaTeX{} is used throughout as an accessibility
aid: its high-quality typesetting, structured navigation,
wheat-coloured page background, and clean separation between source
editing and rendered output all materially reduce visual fatigue
compared to word processors. This is noted here because the same
approach may benefit other learners with similar profiles.

Neuropsychological assessment has identified two features of the
author's cognitive profile that are directly relevant to this
project and worth noting for other learners who may recognise
themselves in them:

\begin{description}

  \item[Reduced working memory] Associated with ADHD, this makes
    rote memorisation of vocabulary lists or paradigm tables
    effortful and largely ineffective as a learning strategy.
    Accordingly, no attempt is made to memorise the material in
    this document. The goal is repeated, pleasurable exposure ---
    the same words and constructions reappearing in new contexts
    until they become familiar without deliberate effort.
    The analyses are written to be enjoyable to read in their own
    right, not to be studied.

  \item[Strong pattern recognition and rule derivation] Identified
    as a compensating strength, this is the capacity to infer
    underlying rules and structures from examples, and to
    generalise them rapidly once the pattern is seen. This is
    precisely what grammatical analysis rewards: rather than
    memorising that \textit{na cmentarzu} means ``at the cemetery,''
    the learner recognises the locative pattern, derives the rule,
    and thereafter reads all locative constructions with growing
    confidence --- without having memorised a single paradigm table.
    The Latin background established this kind of structural
    awareness early; Polish is being acquired in the same way.

\end{description}

These two features together suggest a learning profile that is
poorly served by conventional language instruction (which relies
heavily on working memory and rote practice) but well served by
immersive, contextual, analysis-rich reading of genuinely engaging
material. The method described in this document was developed
organically around exactly this profile.

\section*{The AI Tool}

All phrase-by-phrase translation, grammatical analysis, and literary
commentary in this document was generated interactively using
\textbf{Claude}, an AI assistant developed by
\href{https://www.anthropic.com}{Anthropic}. Specifically, the model
used was \textbf{Claude Sonnet 4.6}
(model string: \texttt{claude-sonnet-4-6}), accessed via the
\href{https://claude.ai}{Claude.ai} web interface.

The \LaTeX{} document itself --- its structure, preamble, custom
commands, grammar reference, and chapter files --- was also drafted
and maintained by the model in response to the author's direction,
updated periodically as new material accumulated.

\section*{The Method}

The workflow is simple and requires no technical expertise beyond
basic \LaTeX{} compilation:

\begin{enumerate}

  \item The learner opens the source book in an e-reader application.
  Since the book is protected by digital rights management (DRM),
  text cannot be copied directly. Instead, the learner takes a
  screenshot of one sentence or short passage at a time.

  \item The screenshot is pasted into a conversation with Claude.
  The model reads the Polish text from the image and provides:
  \begin{itemize}
    \item a full English translation of the passage;
    \item a word-by-word or phrase-by-phrase breakdown with
      vocabulary glosses;
    \item grammatical analysis --- case endings, verb aspect,
      participial constructions, and other relevant features;
    \item literary and contextual commentary where the author's
      choices are particularly interesting.
  \end{itemize}

  \item The learner reads the analysis, asks follow-up questions
  (about vocabulary, grammar, or usage), and continues to the
  next passage at their own pace. There is no pressure to
  memorise --- the goal is absorption through engaged reading.

  \item Periodically, at natural section breaks, the accumulated
  analyses are consolidated into the \LaTeX{} document. The model
  drafts the \LaTeX{} source; the learner compiles it locally
  and reviews the output.

  \item When the AI session approaches its context window limit,
  a new session is started. The main \texttt{.tex} file --- which
  contains a detailed comment block describing the project,
  learner profile, and working method --- is uploaded to the new
  session as a handoff document, allowing work to resume
  seamlessly.

\end{enumerate}

\section*{Depth of Analysis}

The depth of analysis is intentionally variable. Some sentences
are unpacked in full detail because the grammar or vocabulary is
instructive. Others are translated more briefly because the content
is transitional or the constructions are already familiar. The
learner directs the pace: tangential questions about grammar or
usage are encouraged at any point, and the conversation moves
between close analysis and broader literary observation as the
text warrants.

This variability is a feature, not a deficiency. It mirrors the
way a knowledgeable reading companion would engage with a text ---
pausing where things are interesting, moving on where they are not.

\section*{Suitability for Other Languages and Learners}

The method is not specific to Polish. It is applicable to any
language for which a sufficiently capable LLM exists, and to any
learner who has a genuine motivation to read a specific text in
that language. The key conditions that make it work are:

\begin{itemize}
  \item a source text that the learner \textit{actually wants to read}
    --- ideally one unavailable in their native language, removing
    the option of simply reading a translation;
  \item willingness to proceed slowly and to treat the analysis as
    enjoyable in its own right, not merely as a means to an end;
  \item basic familiarity with \LaTeX{} if the learner wishes to
    maintain a compiled document, though the conversational
    analysis alone is valuable without any document production.
\end{itemize}

The method works particularly well for learners with existing
metalinguistic awareness --- those who have studied another
inflected language (Latin, Russian, German, classical Greek) and
already understand concepts like case, aspect, and agreement
abstractly. For such learners, the AI analysis illuminates the
\textit{Polish instantiation} of familiar concepts rather than
introducing entirely new ideas from scratch.

\section*{Transparency and Limitations}

The analyses in this document reflect the capabilities of the model
at the time of writing and should be treated as a knowledgeable
first-pass rather than authoritative scholarship. Occasional errors
in grammatical analysis or nuance are possible. Where the author
has Polish-speaking family or other expert resources available,
those remain the final arbiter of correctness.

The quoted Polish passages in this document are brief and used
solely for educational commentary. They do not constitute
reproduction of the source text and are consistent with fair use
principles for personal, non-commercial educational purposes.
The document is shared openly but the source book remains under
its original copyright.

\vspace{1.5em}
\noindent The full project, including all \LaTeX{} source files,
is available at:\\
\href{https://github.com/}{[GitHub repository URL to be added]}

\vspace{0.5em}
\noindent Readers are welcome to adapt the structure, commands,
and approach for their own language learning projects.

% ============================================================

\chapter*{A Note on Method: Use of Artificial Intelligence}
\addcontentsline{toc}{chapter}{A Note on Method: Use of AI}
\markboth{A Note on Method}{}

This document was produced with the assistance of a large language
model (LLM) and is shared openly in the hope that others may adapt
the method for their own language learning.

\section*{The Learning Context}

The author is a native English speaker with a Polish wife, Polish
in-laws, and a recently acquired apartment in Kraków. After more than
twenty years of incidental exposure to Polish, conversational ability
had plateaued at a level sufficient for simple transactions but
inadequate for real two-way communication. The goal of this project
is to break through that plateau by engaging deeply with a Polish-language
book that is genuinely interesting, not available in English translation,
and whose subject matter --- a lead poisoning epidemic in industrial
Silesia --- connects directly to the author's professional background
in biomedical engineering.

The author has ADHD and dyslexia/dysorthographia, and shares this
openly in the hope of reducing the stigma that still surrounds these
conditions --- particularly in academic and language-learning contexts,
where they are often misread as indicators of low ability or
insufficient effort. \LaTeX{} is used throughout as an accessibility
aid: its high-quality typesetting, structured navigation,
wheat-coloured page background, and clean separation between source
editing and rendered output all materially reduce visual fatigue
compared to word processors. This is noted here because the same
approach may benefit other learners with similar profiles.

Neuropsychological assessment has identified two features of the
author's cognitive profile that are directly relevant to this
project and worth noting for other learners who may recognise
themselves in them:

\begin{description}

  \item[Reduced working memory] Associated with ADHD, this makes
    rote memorisation of vocabulary lists or paradigm tables
    effortful and largely ineffective as a learning strategy.
    Accordingly, no attempt is made to memorise the material in
    this document. The goal is repeated, pleasurable exposure ---
    the same words and constructions reappearing in new contexts
    until they become familiar without deliberate effort.
    The analyses are written to be enjoyable to read in their own
    right, not to be studied.

  \item[Strong pattern recognition and rule derivation] Identified
    as a compensating strength, this is the capacity to infer
    underlying rules and structures from examples, and to
    generalise them rapidly once the pattern is seen. This is
    precisely what grammatical analysis rewards: rather than
    memorising that \textit{na cmentarzu} means ``at the cemetery,''
    the learner recognises the locative pattern, derives the rule,
    and thereafter reads all locative constructions with growing
    confidence --- without having memorised a single paradigm table.
    The Latin background established this kind of structural
    awareness early; Polish is being acquired in the same way.

\end{description}

These two features together suggest a learning profile that is
poorly served by conventional language instruction (which relies
heavily on working memory and rote practice) but well served by
immersive, contextual, analysis-rich reading of genuinely engaging
material. The method described in this document was developed
organically around exactly this profile.

\section*{The AI Tool}

All phrase-by-phrase translation, grammatical analysis, and literary
commentary in this document was generated interactively using
\textbf{Claude}, an AI assistant developed by
\href{https://www.anthropic.com}{Anthropic}. Specifically, the model
used was \textbf{Claude Sonnet 4.6}
(model string: \texttt{claude-sonnet-4-6}), accessed via the
\href{https://claude.ai}{Claude.ai} web interface.

The \LaTeX{} document itself --- its structure, preamble, custom
commands, grammar reference, and chapter files --- was also drafted
and maintained by the model in response to the author's direction,
updated periodically as new material accumulated.

\section*{The Method}

The workflow is simple and requires no technical expertise beyond
basic \LaTeX{} compilation:

\begin{enumerate}

  \item The learner opens the source book in an e-reader application.
  Since the book is protected by digital rights management (DRM),
  text cannot be copied directly. Instead, the learner takes a
  screenshot of one sentence or short passage at a time.

  \item The screenshot is pasted into a conversation with Claude.
  The model reads the Polish text from the image and provides:
  \begin{itemize}
    \item a full English translation of the passage;
    \item a word-by-word or phrase-by-phrase breakdown with
      vocabulary glosses;
    \item grammatical analysis --- case endings, verb aspect,
      participial constructions, and other relevant features;
    \item literary and contextual commentary where the author's
      choices are particularly interesting.
  \end{itemize}

  \item The learner reads the analysis, asks follow-up questions
  (about vocabulary, grammar, or usage), and continues to the
  next passage at their own pace. There is no pressure to
  memorise --- the goal is absorption through engaged reading.

  \item Periodically, at natural section breaks, the accumulated
  analyses are consolidated into the \LaTeX{} document. The model
  drafts the \LaTeX{} source; the learner compiles it locally
  and reviews the output.

  \item When the AI session approaches its context window limit,
  a new session is started. The main \texttt{.tex} file --- which
  contains a detailed comment block describing the project,
  learner profile, and working method --- is uploaded to the new
  session as a handoff document, allowing work to resume
  seamlessly.

\end{enumerate}

\section*{Depth of Analysis}

The depth of analysis is intentionally variable. Some sentences
are unpacked in full detail because the grammar or vocabulary is
instructive. Others are translated more briefly because the content
is transitional or the constructions are already familiar. The
learner directs the pace: tangential questions about grammar or
usage are encouraged at any point, and the conversation moves
between close analysis and broader literary observation as the
text warrants.

This variability is a feature, not a deficiency. It mirrors the
way a knowledgeable reading companion would engage with a text ---
pausing where things are interesting, moving on where they are not.

\section*{Suitability for Other Languages and Learners}

The method is not specific to Polish. It is applicable to any
language for which a sufficiently capable LLM exists, and to any
learner who has a genuine motivation to read a specific text in
that language. The key conditions that make it work are:

\begin{itemize}
  \item a source text that the learner \textit{actually wants to read}
    --- ideally one unavailable in their native language, removing
    the option of simply reading a translation;
  \item willingness to proceed slowly and to treat the analysis as
    enjoyable in its own right, not merely as a means to an end;
  \item basic familiarity with \LaTeX{} if the learner wishes to
    maintain a compiled document, though the conversational
    analysis alone is valuable without any document production.
\end{itemize}

The method works particularly well for learners with existing
metalinguistic awareness --- those who have studied another
inflected language (Latin, Russian, German, classical Greek) and
already understand concepts like case, aspect, and agreement
abstractly. For such learners, the AI analysis illuminates the
\textit{Polish instantiation} of familiar concepts rather than
introducing entirely new ideas from scratch.

\section*{Transparency and Limitations}

The analyses in this document reflect the capabilities of the model
at the time of writing and should be treated as a knowledgeable
first-pass rather than authoritative scholarship. Occasional errors
in grammatical analysis or nuance are possible. Where the author
has Polish-speaking family or other expert resources available,
those remain the final arbiter of correctness.

The quoted Polish passages in this document are brief and used
solely for educational commentary. They do not constitute
reproduction of the source text and are consistent with fair use
principles for personal, non-commercial educational purposes.
The document is shared openly but the source book remains under
its original copyright.

\vspace{1.5em}
\noindent The full project, including all \LaTeX{} source files,
is available at:\\
\href{https://github.com/}{[GitHub repository URL to be added]}

\vspace{0.5em}
\noindent Readers are welcome to adapt the structure, commands,
and approach for their own language learning projects.

% ============================================================

\chapter*{A Note on Method: Use of Artificial Intelligence}
\addcontentsline{toc}{chapter}{A Note on Method: Use of AI}
\markboth{A Note on Method}{}

This document was produced with the assistance of a large language
model (LLM) and is shared openly in the hope that others may adapt
the method for their own language learning.

\section*{The Learning Context}

The author is a native English speaker with a Polish wife, Polish
in-laws, and a recently acquired apartment in Kraków. After more than
twenty years of incidental exposure to Polish, conversational ability
had plateaued at a level sufficient for simple transactions but
inadequate for real two-way communication. The goal of this project
is to break through that plateau by engaging deeply with a Polish-language
book that is genuinely interesting, not available in English translation,
and whose subject matter --- a lead poisoning epidemic in industrial
Silesia --- connects directly to the author's professional background
in biomedical engineering.

The author has ADHD and dyslexia/dysorthographia, and shares this
openly in the hope of reducing the stigma that still surrounds these
conditions --- particularly in academic and language-learning contexts,
where they are often misread as indicators of low ability or
insufficient effort. \LaTeX{} is used throughout as an accessibility
aid: its high-quality typesetting, structured navigation,
wheat-coloured page background, and clean separation between source
editing and rendered output all materially reduce visual fatigue
compared to word processors. This is noted here because the same
approach may benefit other learners with similar profiles.

Neuropsychological assessment has identified two features of the
author's cognitive profile that are directly relevant to this
project and worth noting for other learners who may recognise
themselves in them:

\begin{description}

  \item[Reduced working memory] Associated with ADHD, this makes
    rote memorisation of vocabulary lists or paradigm tables
    effortful and largely ineffective as a learning strategy.
    Accordingly, no attempt is made to memorise the material in
    this document. The goal is repeated, pleasurable exposure ---
    the same words and constructions reappearing in new contexts
    until they become familiar without deliberate effort.
    The analyses are written to be enjoyable to read in their own
    right, not to be studied.

  \item[Strong pattern recognition and rule derivation] Identified
    as a compensating strength, this is the capacity to infer
    underlying rules and structures from examples, and to
    generalise them rapidly once the pattern is seen. This is
    precisely what grammatical analysis rewards: rather than
    memorising that \textit{na cmentarzu} means ``at the cemetery,''
    the learner recognises the locative pattern, derives the rule,
    and thereafter reads all locative constructions with growing
    confidence --- without having memorised a single paradigm table.
    The Latin background established this kind of structural
    awareness early; Polish is being acquired in the same way.

\end{description}

These two features together suggest a learning profile that is
poorly served by conventional language instruction (which relies
heavily on working memory and rote practice) but well served by
immersive, contextual, analysis-rich reading of genuinely engaging
material. The method described in this document was developed
organically around exactly this profile.

\section*{The AI Tool}

All phrase-by-phrase translation, grammatical analysis, and literary
commentary in this document was generated interactively using
\textbf{Claude}, an AI assistant developed by
\href{https://www.anthropic.com}{Anthropic}. Specifically, the model
used was \textbf{Claude Sonnet 4.6}
(model string: \texttt{claude-sonnet-4-6}), accessed via the
\href{https://claude.ai}{Claude.ai} web interface.

The \LaTeX{} document itself --- its structure, preamble, custom
commands, grammar reference, and chapter files --- was also drafted
and maintained by the model in response to the author's direction,
updated periodically as new material accumulated.

\section*{The Method}

The workflow is simple and requires no technical expertise beyond
basic \LaTeX{} compilation:

\begin{enumerate}

  \item The learner opens the source book in an e-reader application.
  Since the book is protected by digital rights management (DRM),
  text cannot be copied directly. Instead, the learner takes a
  screenshot of one sentence or short passage at a time.

  \item The screenshot is pasted into a conversation with Claude.
  The model reads the Polish text from the image and provides:
  \begin{itemize}
    \item a full English translation of the passage;
    \item a word-by-word or phrase-by-phrase breakdown with
      vocabulary glosses;
    \item grammatical analysis --- case endings, verb aspect,
      participial constructions, and other relevant features;
    \item literary and contextual commentary where the author's
      choices are particularly interesting.
  \end{itemize}

  \item The learner reads the analysis, asks follow-up questions
  (about vocabulary, grammar, or usage), and continues to the
  next passage at their own pace. There is no pressure to
  memorise --- the goal is absorption through engaged reading.

  \item Periodically, at natural section breaks, the accumulated
  analyses are consolidated into the \LaTeX{} document. The model
  drafts the \LaTeX{} source; the learner compiles it locally
  and reviews the output.

  \item When the AI session approaches its context window limit,
  a new session is started. The main \texttt{.tex} file --- which
  contains a detailed comment block describing the project,
  learner profile, and working method --- is uploaded to the new
  session as a handoff document, allowing work to resume
  seamlessly.

\end{enumerate}

\section*{Depth of Analysis}

The depth of analysis is intentionally variable. Some sentences
are unpacked in full detail because the grammar or vocabulary is
instructive. Others are translated more briefly because the content
is transitional or the constructions are already familiar. The
learner directs the pace: tangential questions about grammar or
usage are encouraged at any point, and the conversation moves
between close analysis and broader literary observation as the
text warrants.

This variability is a feature, not a deficiency. It mirrors the
way a knowledgeable reading companion would engage with a text ---
pausing where things are interesting, moving on where they are not.

\section*{Suitability for Other Languages and Learners}

The method is not specific to Polish. It is applicable to any
language for which a sufficiently capable LLM exists, and to any
learner who has a genuine motivation to read a specific text in
that language. The key conditions that make it work are:

\begin{itemize}
  \item a source text that the learner \textit{actually wants to read}
    --- ideally one unavailable in their native language, removing
    the option of simply reading a translation;
  \item willingness to proceed slowly and to treat the analysis as
    enjoyable in its own right, not merely as a means to an end;
  \item basic familiarity with \LaTeX{} if the learner wishes to
    maintain a compiled document, though the conversational
    analysis alone is valuable without any document production.
\end{itemize}

The method works particularly well for learners with existing
metalinguistic awareness --- those who have studied another
inflected language (Latin, Russian, German, classical Greek) and
already understand concepts like case, aspect, and agreement
abstractly. For such learners, the AI analysis illuminates the
\textit{Polish instantiation} of familiar concepts rather than
introducing entirely new ideas from scratch.

\section*{Transparency and Limitations}

The analyses in this document reflect the capabilities of the model
at the time of writing and should be treated as a knowledgeable
first-pass rather than authoritative scholarship. Occasional errors
in grammatical analysis or nuance are possible. Where the author
has Polish-speaking family or other expert resources available,
those remain the final arbiter of correctness.

The quoted Polish passages in this document are brief and used
solely for educational commentary. They do not constitute
reproduction of the source text and are consistent with fair use
principles for personal, non-commercial educational purposes.
The document is shared openly but the source book remains under
its original copyright.

\vspace{1.5em}
\noindent The full project, including all \LaTeX{} source files,
is available at:\\
\href{https://github.com/}{[GitHub repository URL to be added]}

\vspace{0.5em}
\noindent Readers are welcome to adapt the structure, commands,
and approach for their own language learning projects.
}
  {% ============================================================
%  preface_ai_short.tex  —  Abbreviated AI disclosure (< half page)
%  For the full account of method, see preface_ai.tex
% ============================================================

\chapter*{A Note on Method}
\addcontentsline{toc}{chapter}{A Note on Method}
\markboth{A Note on Method}{}

The phrase-by-phrase analyses in this document were produced
interactively with the assistance of Anthropic's Claude (model:
\texttt{claude-sonnet-4-6}), accessed via the Claude.ai web interface.
Working from screenshots of the source text supplied one passage at a
time, the model provided translation, word-level grammatical breakdown,
and literary commentary in a conversational session. The \LaTeX{}
document was also structured and updated by the model throughout.

Quoted passages are brief and used solely for educational commentary.
This document is shared openly --- source code and compiled PDF are
available in a public repository on GitHub:

\begin{center}
\url{https://github.com/SnowdonAnalytica/olowiane-dzieci-analiza}
\end{center}

A full account of the method --- including the learner profile,
rationale, and workflow --- is available in the companion file
\texttt{preface\_ai.tex}.
}

% ── Grammar Foundations ──────────────────────────────────────
% ============================================================
%  grammar.tex  —  Polish Grammar Foundations
%  Referenced by polish_study.tex via % ============================================================
%  grammar.tex  —  Polish Grammar Foundations
%  Referenced by polish_study.tex via % ============================================================
%  grammar.tex  —  Polish Grammar Foundations
%  Referenced by polish_study.tex via \input{grammar}
% ============================================================

\chapter*{Polish Grammar Foundations}
\addcontentsline{toc}{chapter}{Polish Grammar Foundations}
\markboth{Polish Grammar Foundations}{}

This chapter provides a compact reference for the grammatical structures 
encountered throughout the reading. It is not intended as a complete grammar 
— rather, as a touchstone for constructs that recur in the text. Cross-references 
to passages where each construct first appears are noted where relevant.

% ────────────────────────────────────────────────────────────
\section{Nouns and Gender}

Polish nouns have three grammatical genders: \textbf{masculine}, \textbf{feminine}, 
and \textbf{neuter}. Unlike French, where gender must largely be memorised, Polish 
gender is often inferable from the noun's ending.

\begin{description}
  \item[Masculine] Typically end in a consonant: \polword{grób}{grave}, 
    \polword{pogrzeb}{funeral}, \polword{syn}{son}.
  \item[Feminine] Typically end in \textbf{-a} or \textbf{-i/y}: 
    \polword{wdowa}{widow}, \polword{grupka}{small group}, 
    \polword{nadzieja}{hope}.
  \item[Neuter] Typically end in \textbf{-o} or \textbf{-e}: 
    \polword{dzieciństwo}{childhood}, \polword{dziecko}{child}.
\end{description}

\grambox{Note on animacy}{Masculine nouns are further divided into 
\textit{animate} (persons and animals) and \textit{inanimate}. This distinction 
affects the accusative case ending and is particularly important for direct objects.}

Nouns also form diminutives extensively in Polish, often with suffixes 
\textbf{-ka}, \textbf{-ko}, \textbf{-ek}, \textbf{-eczka}. These indicate 
smallness but frequently carry emotional colouring — affection, sympathy, or 
irony — rather than literal size. Examples encountered in Chapter 1: 
\polword{grupka}{little group} (from \textit{grupa}), 
\polword{garstka}{small handful} (from \textit{garść}), 
\polword{flaszka}{bottle, colloquial} (from \textit{flasza}),
\polword{podstawówka}{primary school, affectionate} (from \textit{szkoła podstawowa}).

% ────────────────────────────────────────────────────────────
\section{The Seven Cases}

Polish has seven grammatical cases. Each case encodes a grammatical 
relationship that English expresses through word order or prepositions. 
The Latin equivalents are noted as a reference point.

\begin{description}

  \item[Nominative (\textit{mianownik})] The subject of the sentence. 
    Who or what performs the action. Equivalent to Latin nominative.\\
    \textit{Groszek \textbf{powiedział}\ldots} — Groszek said\ldots

  \item[Genitive (\textit{dopełniacz})] Possession, absence, quantity, 
    and the object of many prepositions (\textit{do, od, z, bez, dla, oprócz}). 
    The most frequently used case after the nominative. Equivalent to Latin genitive.\\
    \textit{dla \textbf{uczczenia} jego \textbf{pamięci}} — for the honouring of his memory\\
    \textit{garstki \textbf{obecnych}} — a handful of those present

  \item[Dative (\textit{celownik})] The indirect object — to whom or for whom. 
    Equivalent to Latin dative. Encodes dedication, as seen in the book's 
    epigraph: \textit{\textbf{Ojcu}, \textbf{Synowi}} — to [my] father, to [my] son.

  \item[Accusative (\textit{biernik})] The direct object of most verbs, 
    and the object of motion prepositions (\textit{na, w, przez, za}).
    Equivalent to Latin accusative.\\
    \textit{Zostaw \textbf{tę flaszkę}!} — Leave that bottle!

  \item[Instrumental (\textit{narzędnik})] Means, accompaniment, and 
    identity (used after the verb \textit{być} — to be, for professions 
    and states). Also follows the preposition \textit{z} (with) and \textit{za} 
    (behind/after in positional sense). Equivalent to Latin ablative of means.\\
    \textit{z \textbf{psem}} — with a dog\\
    \textit{z \textbf{wykształcenia} inżynier} — an engineer by training

  \item[Locative (\textit{miejscownik})] Location and topic. 
    \textbf{Always} used with a preposition (\textit{w, na, przy, o, po}).
    Never appears without one. Equivalent to Latin ablative of place.\\
    \textit{na \textbf{cmentarzu}} — at the cemetery\\
    \textit{przy \textbf{grobie}} — by the grave

  \item[Vocative (\textit{wołacz})] Direct address. Used when calling 
    or speaking to someone by name or title. Relatively rare in prose 
    but common in dialogue.\\
    \textit{Mamo!} — Mum! (addressing one's mother)

\end{description}

\grambox{Key principle}{Case endings change for both nouns \textit{and} 
adjectives — they must always agree in gender, number, and case. 
This is identical to Latin adjective-noun concord: 
\textit{wczesną jesienią} (in early autumn) — both words take the 
instrumental feminine singular ending.}

% ────────────────────────────────────────────────────────────
\section{Adjectives}

Adjectives agree with their nouns in gender, number, and case. The 
ending of the adjective therefore changes to match the noun it modifies, 
across all seven cases and three genders — 42 potential forms in principle, 
though many are identical in practice.

\textit{świeżo wykopany grób} (freshly dug grave, masc. nom.) — the participle 
\textit{wykopany} agrees with \textit{grób}.

\textit{świeżo wykopanym grobie} (by the freshly dug grave, masc. loc.) — both 
adjective and noun shift to the locative.

% ────────────────────────────────────────────────────────────
\section{Verbs: Aspect}

The single most important concept in Polish verbal grammar, and the one 
with no direct Latin equivalent, is \textbf{verbal aspect}.

Every Polish verb exists in two versions:

\begin{description}
  \item[Imperfective (\textit{niedokonany})] Describes ongoing, repeated, 
    or incomplete actions. What was happening, what happens habitually, 
    what is in progress.
  \item[Perfective (\textit{dokonany})] Describes completed, bounded, 
    or one-time actions. What happened (and is done), what will definitively happen.
\end{description}

\grambox{Key contrast}{\polword{pisać}{to write (imperfective, ongoing)} 
vs.\ \polword{napisać}{to write (perfective, completed)}. The prefix 
\textit{na-} here creates the perfective. This pattern — imperfective root, 
prefixed perfective — is one of the most common in Polish.}

Aspect affects \textbf{what tenses are possible}. Perfective verbs have 
no present tense (a completed action cannot be in progress now). Their 
``present-looking'' form is actually future: \textit{napiszę} = I will write (and finish it).

\subsection*{Common perfectivising prefixes}

These prefixes convert imperfective verbs to perfective and are extremely 
productive in Polish. Recognising them unlocks large swaths of vocabulary.

\begin{description}
  \item[\textbf{za-}] Often signals the onset or completion of an action: 
    \polword{pamiętać}{to remember} → \polword{zapomnieć}{to forget}; 
    \polword{truć}{to poison} → \polword{zatruć}{to poison (completely)}.
  \item[\textbf{po-}] Often signals completion or doing something for a while: 
    \polword{czekać}{to wait} → \polword{poczekać}{to wait (for a bit / done waiting)}.
  \item[\textbf{na-}] Completion, sufficiency: 
    \polword{pisać}{to write} → \polword{napisać}{to write (and finish)}.
  \item[\textbf{od-}] Away from, completing a departure: 
    \polword{iść}{to go} → \polword{odejść}{to go away, depart}.
  \item[\textbf{wy-}] Out, thoroughly: 
    \polword{kopać}{to dig} → \polword{wykopać}{to dig out/up}.
  \item[\textbf{z- / s-}] Completion, bringing together: 
    \polword{bliski}{close} → \polword{zbliżyć się}{to draw closer}.
\end{description}

% ────────────────────────────────────────────────────────────
\section{Verb Conjugation}

Polish verbs encode person and number in their endings, making subject 
pronouns optional (as in Latin and Russian). The pronouns are used 
for emphasis or contrast.

\subsection*{Present tense — example: \textit{czekać} (to wait, imperfective)}

\begin{tabular}{lll}
  \textbf{Person} & \textbf{Singular} & \textbf{Plural} \\
  1st & czekam (I wait) & czekamy (we wait) \\
  2nd & czekasz (you wait) & czekacie (you wait) \\
  3rd & czeka (he/she/it waits) & czekają (they wait) \\
\end{tabular}

\subsection*{Past tense — gender marking}

Uniquely among European languages, Polish past tense verbs agree with 
the gender of the subject. This gender marking is a suffix added to the 
past stem:

\begin{description}
  \item[Masculine singular] \textbf{-ł}: \textit{powiedział} (he said), 
    \textit{spotkałem} (I [masc.] met)
  \item[Feminine singular] \textbf{-ła}: \textit{czekała} (she waited), 
    \textit{musiała} (she had to)
  \item[Neuter singular] \textbf{-ło}: \textit{miało} (it was supposed to)
  \item[Virile plural (masc. personal)] \textbf{-li}: \textit{zdecydowaliśmy} 
    (we [men/mixed] decided), \textit{ustawili się} (they positioned themselves)
  \item[Non-virile plural] \textbf{-ły}: used for groups of women or non-persons
\end{description}

\grambox{Virile vs.\ non-virile}{Polish distinguishes between groups 
containing at least one man (\textit{virile}, personal masculine plural) 
and all other plurals. This distinction affects verb endings, adjective 
forms, and pronouns throughout the grammar.}

% ────────────────────────────────────────────────────────────
\section{Key Impersonal Constructions}

Polish uses several impersonal constructions that have no direct English 
equivalent but will be encountered on virtually every page.

\begin{description}
  \item[\polword{trzeba}{one must, it is necessary}] Followed by an 
    infinitive. No subject. \textit{Trzeba kontynuować} — one must continue. 
    Equivalent to French \textit{il faut}.
  \item[\polword{można}{one can, it is possible}] Followed by an 
    infinitive. \textit{Można było tam dostać} — one could get there. 
    Equivalent to French \textit{on peut}.
  \item[\polword{warto}{it is worth}] Followed by an infinitive. 
    \textit{Warto wiedzieć} — it is worth knowing.
\end{description}

% ────────────────────────────────────────────────────────────
\section{Essential Particles and Connectives}

A small set of particles carries enormous communicative weight in Polish 
and will be encountered constantly.

\begin{description}
  \item[\polword{już}{already, now, anymore}] Marks a transition or 
    change of state. \textit{Już rozumiem} — I understand now (I've crossed 
    into understanding). With negation: \textit{już nie} — not anymore.
  \item[\polword{jeszcze}{still, yet, even more}] Marks continuation 
    or addition. \textit{Jeszcze jeden} — yet one more. With negation: 
    \textit{jeszcze nie} — not yet.
  \item[\polword{już nie}{not anymore}] Something was happening but has stopped.
  \item[\polword{jeszcze nie}{not yet}] Something has not happened but may.
  \item[\polword{też / także}{also, too}] \textit{Też} is conversational; 
    \textit{także} is slightly more formal.
  \item[\polword{wcale nie}{not at all}] Emphatic negation. 
    \textit{Wcale nie odchodzi} — it does not go away at all.
  \item[\polword{właśnie}{exactly, just now, precisely}] 
    Confirms or emphasises: \textit{właśnie} as a standalone response 
    means ``exactly'' or ``quite so.''
  \item[\polword{chyba}{probably, I suppose, surely}] Expresses 
    uncertainty or mild assertion. Very common in speech.
  \item[\polword{przecież}{after all, but surely}] Implies that 
    something should be obvious. \textit{Przecież wiesz} — you know 
    this perfectly well (after all).
\end{description}

% ────────────────────────────────────────────────────────────
\section{Conjunctions: Because}

Three words translate as ``because'' with different registers:

\begin{description}
  \item[\polword{bo}{because}] Casual, spoken, everyday. 
    The default in conversation.
  \item[\polword{ponieważ}{because}] Standard written and formal spoken. 
    Neutral and widely applicable.
  \item[\polword{gdyż}{because/for}] Literary and elevated. 
    Signals deliberate, crafted prose. Equivalent to French \textit{car}.
\end{description}

% ────────────────────────────────────────────────────────────
\section{Word Order}

Polish word order is considerably freer than English, because case 
endings — not position — identify grammatical role. However, order 
is far from arbitrary: it encodes \textbf{information structure}.

\begin{description}
  \item[New information] tends to come at the \textbf{end} of the sentence 
    (the ``rheme'' or focus position). This is why Polish often places the 
    subject after the verb, or saves the most concrete noun for last.
  \item[Given / known information] tends to come early.
  \item[Verb-first] is used for dramatic effect, as in: 
    \textit{Umarł Fiszer} — Fiszer died (the death before the name).
\end{description}

This freedom means that close translation often fails to capture 
the author's emphasis. Notes in the chapter analyses point out 
significant word-order choices where they carry meaning.

% ────────────────────────────────────────────────────────────
\section{Reflexive Verbs and \textit{się}}

The particle \textbf{się} is one of the most versatile elements in 
Polish. It functions as:

\begin{description}
  \item[Reflexive marker] The action reflects back on the subject: 
    \textit{zbliżyć się} — to bring oneself closer, to approach.
  \item[Reciprocal marker] \textit{spotkać się} — to meet (each other).
  \item[Impersonal/passive marker] \textit{mówi się} — one says, 
    it is said. Removes the agent entirely.
  \item[Part of the verb's inherent meaning] Some verbs simply 
    require \textit{się} with no separable reflexive meaning: 
    \textit{bać się} (to be afraid), \textit{cieszyć się} (to be glad).
\end{description}

% ────────────────────────────────────────────────────────────
\section{Participles}

Polish uses participial constructions extensively — more so than 
conversational English, but comparable to written French or Latin.

\begin{description}
  \item[Present active participle (\textit{imiesłów przymiotnikowy czynny})] 
    Describes an ongoing action of the noun it modifies. 
    Formed with suffix \textbf{-ący/-ąca/-ące}: 
    \textit{grupka \textbf{stojąca} przy grobie} — the group standing by the grave.
  \item[Past passive participle (\textit{imiesłów przymiotnikowy bierny})] 
    Describes a completed action done to the noun. 
    Formed with suffix \textbf{-ny/-na/-ne} or \textbf{-ty/-ta/-te}: 
    \textit{świeżo \textbf{wykopany} grób} — the freshly dug grave; 
    \textit{zapomniana epidemia} — the forgotten epidemic.
  \item[Adverbial participle (\textit{imiesłów przysłówkowy})] 
    Describes a simultaneous action of the subject. Does not decline. 
    Formed with suffix \textbf{-ąc}: 
    \textit{\textbf{odkręcając} nakrętkę} — while unscrewing the cap.
\end{description}

% ============================================================

\chapter*{Polish Grammar Foundations}
\addcontentsline{toc}{chapter}{Polish Grammar Foundations}
\markboth{Polish Grammar Foundations}{}

This chapter provides a compact reference for the grammatical structures 
encountered throughout the reading. It is not intended as a complete grammar 
— rather, as a touchstone for constructs that recur in the text. Cross-references 
to passages where each construct first appears are noted where relevant.

% ────────────────────────────────────────────────────────────
\section{Nouns and Gender}

Polish nouns have three grammatical genders: \textbf{masculine}, \textbf{feminine}, 
and \textbf{neuter}. Unlike French, where gender must largely be memorised, Polish 
gender is often inferable from the noun's ending.

\begin{description}
  \item[Masculine] Typically end in a consonant: \polword{grób}{grave}, 
    \polword{pogrzeb}{funeral}, \polword{syn}{son}.
  \item[Feminine] Typically end in \textbf{-a} or \textbf{-i/y}: 
    \polword{wdowa}{widow}, \polword{grupka}{small group}, 
    \polword{nadzieja}{hope}.
  \item[Neuter] Typically end in \textbf{-o} or \textbf{-e}: 
    \polword{dzieciństwo}{childhood}, \polword{dziecko}{child}.
\end{description}

\grambox{Note on animacy}{Masculine nouns are further divided into 
\textit{animate} (persons and animals) and \textit{inanimate}. This distinction 
affects the accusative case ending and is particularly important for direct objects.}

Nouns also form diminutives extensively in Polish, often with suffixes 
\textbf{-ka}, \textbf{-ko}, \textbf{-ek}, \textbf{-eczka}. These indicate 
smallness but frequently carry emotional colouring — affection, sympathy, or 
irony — rather than literal size. Examples encountered in Chapter 1: 
\polword{grupka}{little group} (from \textit{grupa}), 
\polword{garstka}{small handful} (from \textit{garść}), 
\polword{flaszka}{bottle, colloquial} (from \textit{flasza}),
\polword{podstawówka}{primary school, affectionate} (from \textit{szkoła podstawowa}).

% ────────────────────────────────────────────────────────────
\section{The Seven Cases}

Polish has seven grammatical cases. Each case encodes a grammatical 
relationship that English expresses through word order or prepositions. 
The Latin equivalents are noted as a reference point.

\begin{description}

  \item[Nominative (\textit{mianownik})] The subject of the sentence. 
    Who or what performs the action. Equivalent to Latin nominative.\\
    \textit{Groszek \textbf{powiedział}\ldots} — Groszek said\ldots

  \item[Genitive (\textit{dopełniacz})] Possession, absence, quantity, 
    and the object of many prepositions (\textit{do, od, z, bez, dla, oprócz}). 
    The most frequently used case after the nominative. Equivalent to Latin genitive.\\
    \textit{dla \textbf{uczczenia} jego \textbf{pamięci}} — for the honouring of his memory\\
    \textit{garstki \textbf{obecnych}} — a handful of those present

  \item[Dative (\textit{celownik})] The indirect object — to whom or for whom. 
    Equivalent to Latin dative. Encodes dedication, as seen in the book's 
    epigraph: \textit{\textbf{Ojcu}, \textbf{Synowi}} — to [my] father, to [my] son.

  \item[Accusative (\textit{biernik})] The direct object of most verbs, 
    and the object of motion prepositions (\textit{na, w, przez, za}).
    Equivalent to Latin accusative.\\
    \textit{Zostaw \textbf{tę flaszkę}!} — Leave that bottle!

  \item[Instrumental (\textit{narzędnik})] Means, accompaniment, and 
    identity (used after the verb \textit{być} — to be, for professions 
    and states). Also follows the preposition \textit{z} (with) and \textit{za} 
    (behind/after in positional sense). Equivalent to Latin ablative of means.\\
    \textit{z \textbf{psem}} — with a dog\\
    \textit{z \textbf{wykształcenia} inżynier} — an engineer by training

  \item[Locative (\textit{miejscownik})] Location and topic. 
    \textbf{Always} used with a preposition (\textit{w, na, przy, o, po}).
    Never appears without one. Equivalent to Latin ablative of place.\\
    \textit{na \textbf{cmentarzu}} — at the cemetery\\
    \textit{przy \textbf{grobie}} — by the grave

  \item[Vocative (\textit{wołacz})] Direct address. Used when calling 
    or speaking to someone by name or title. Relatively rare in prose 
    but common in dialogue.\\
    \textit{Mamo!} — Mum! (addressing one's mother)

\end{description}

\grambox{Key principle}{Case endings change for both nouns \textit{and} 
adjectives — they must always agree in gender, number, and case. 
This is identical to Latin adjective-noun concord: 
\textit{wczesną jesienią} (in early autumn) — both words take the 
instrumental feminine singular ending.}

% ────────────────────────────────────────────────────────────
\section{Adjectives}

Adjectives agree with their nouns in gender, number, and case. The 
ending of the adjective therefore changes to match the noun it modifies, 
across all seven cases and three genders — 42 potential forms in principle, 
though many are identical in practice.

\textit{świeżo wykopany grób} (freshly dug grave, masc. nom.) — the participle 
\textit{wykopany} agrees with \textit{grób}.

\textit{świeżo wykopanym grobie} (by the freshly dug grave, masc. loc.) — both 
adjective and noun shift to the locative.

% ────────────────────────────────────────────────────────────
\section{Verbs: Aspect}

The single most important concept in Polish verbal grammar, and the one 
with no direct Latin equivalent, is \textbf{verbal aspect}.

Every Polish verb exists in two versions:

\begin{description}
  \item[Imperfective (\textit{niedokonany})] Describes ongoing, repeated, 
    or incomplete actions. What was happening, what happens habitually, 
    what is in progress.
  \item[Perfective (\textit{dokonany})] Describes completed, bounded, 
    or one-time actions. What happened (and is done), what will definitively happen.
\end{description}

\grambox{Key contrast}{\polword{pisać}{to write (imperfective, ongoing)} 
vs.\ \polword{napisać}{to write (perfective, completed)}. The prefix 
\textit{na-} here creates the perfective. This pattern — imperfective root, 
prefixed perfective — is one of the most common in Polish.}

Aspect affects \textbf{what tenses are possible}. Perfective verbs have 
no present tense (a completed action cannot be in progress now). Their 
``present-looking'' form is actually future: \textit{napiszę} = I will write (and finish it).

\subsection*{Common perfectivising prefixes}

These prefixes convert imperfective verbs to perfective and are extremely 
productive in Polish. Recognising them unlocks large swaths of vocabulary.

\begin{description}
  \item[\textbf{za-}] Often signals the onset or completion of an action: 
    \polword{pamiętać}{to remember} → \polword{zapomnieć}{to forget}; 
    \polword{truć}{to poison} → \polword{zatruć}{to poison (completely)}.
  \item[\textbf{po-}] Often signals completion or doing something for a while: 
    \polword{czekać}{to wait} → \polword{poczekać}{to wait (for a bit / done waiting)}.
  \item[\textbf{na-}] Completion, sufficiency: 
    \polword{pisać}{to write} → \polword{napisać}{to write (and finish)}.
  \item[\textbf{od-}] Away from, completing a departure: 
    \polword{iść}{to go} → \polword{odejść}{to go away, depart}.
  \item[\textbf{wy-}] Out, thoroughly: 
    \polword{kopać}{to dig} → \polword{wykopać}{to dig out/up}.
  \item[\textbf{z- / s-}] Completion, bringing together: 
    \polword{bliski}{close} → \polword{zbliżyć się}{to draw closer}.
\end{description}

% ────────────────────────────────────────────────────────────
\section{Verb Conjugation}

Polish verbs encode person and number in their endings, making subject 
pronouns optional (as in Latin and Russian). The pronouns are used 
for emphasis or contrast.

\subsection*{Present tense — example: \textit{czekać} (to wait, imperfective)}

\begin{tabular}{lll}
  \textbf{Person} & \textbf{Singular} & \textbf{Plural} \\
  1st & czekam (I wait) & czekamy (we wait) \\
  2nd & czekasz (you wait) & czekacie (you wait) \\
  3rd & czeka (he/she/it waits) & czekają (they wait) \\
\end{tabular}

\subsection*{Past tense — gender marking}

Uniquely among European languages, Polish past tense verbs agree with 
the gender of the subject. This gender marking is a suffix added to the 
past stem:

\begin{description}
  \item[Masculine singular] \textbf{-ł}: \textit{powiedział} (he said), 
    \textit{spotkałem} (I [masc.] met)
  \item[Feminine singular] \textbf{-ła}: \textit{czekała} (she waited), 
    \textit{musiała} (she had to)
  \item[Neuter singular] \textbf{-ło}: \textit{miało} (it was supposed to)
  \item[Virile plural (masc. personal)] \textbf{-li}: \textit{zdecydowaliśmy} 
    (we [men/mixed] decided), \textit{ustawili się} (they positioned themselves)
  \item[Non-virile plural] \textbf{-ły}: used for groups of women or non-persons
\end{description}

\grambox{Virile vs.\ non-virile}{Polish distinguishes between groups 
containing at least one man (\textit{virile}, personal masculine plural) 
and all other plurals. This distinction affects verb endings, adjective 
forms, and pronouns throughout the grammar.}

% ────────────────────────────────────────────────────────────
\section{Key Impersonal Constructions}

Polish uses several impersonal constructions that have no direct English 
equivalent but will be encountered on virtually every page.

\begin{description}
  \item[\polword{trzeba}{one must, it is necessary}] Followed by an 
    infinitive. No subject. \textit{Trzeba kontynuować} — one must continue. 
    Equivalent to French \textit{il faut}.
  \item[\polword{można}{one can, it is possible}] Followed by an 
    infinitive. \textit{Można było tam dostać} — one could get there. 
    Equivalent to French \textit{on peut}.
  \item[\polword{warto}{it is worth}] Followed by an infinitive. 
    \textit{Warto wiedzieć} — it is worth knowing.
\end{description}

% ────────────────────────────────────────────────────────────
\section{Essential Particles and Connectives}

A small set of particles carries enormous communicative weight in Polish 
and will be encountered constantly.

\begin{description}
  \item[\polword{już}{already, now, anymore}] Marks a transition or 
    change of state. \textit{Już rozumiem} — I understand now (I've crossed 
    into understanding). With negation: \textit{już nie} — not anymore.
  \item[\polword{jeszcze}{still, yet, even more}] Marks continuation 
    or addition. \textit{Jeszcze jeden} — yet one more. With negation: 
    \textit{jeszcze nie} — not yet.
  \item[\polword{już nie}{not anymore}] Something was happening but has stopped.
  \item[\polword{jeszcze nie}{not yet}] Something has not happened but may.
  \item[\polword{też / także}{also, too}] \textit{Też} is conversational; 
    \textit{także} is slightly more formal.
  \item[\polword{wcale nie}{not at all}] Emphatic negation. 
    \textit{Wcale nie odchodzi} — it does not go away at all.
  \item[\polword{właśnie}{exactly, just now, precisely}] 
    Confirms or emphasises: \textit{właśnie} as a standalone response 
    means ``exactly'' or ``quite so.''
  \item[\polword{chyba}{probably, I suppose, surely}] Expresses 
    uncertainty or mild assertion. Very common in speech.
  \item[\polword{przecież}{after all, but surely}] Implies that 
    something should be obvious. \textit{Przecież wiesz} — you know 
    this perfectly well (after all).
\end{description}

% ────────────────────────────────────────────────────────────
\section{Conjunctions: Because}

Three words translate as ``because'' with different registers:

\begin{description}
  \item[\polword{bo}{because}] Casual, spoken, everyday. 
    The default in conversation.
  \item[\polword{ponieważ}{because}] Standard written and formal spoken. 
    Neutral and widely applicable.
  \item[\polword{gdyż}{because/for}] Literary and elevated. 
    Signals deliberate, crafted prose. Equivalent to French \textit{car}.
\end{description}

% ────────────────────────────────────────────────────────────
\section{Word Order}

Polish word order is considerably freer than English, because case 
endings — not position — identify grammatical role. However, order 
is far from arbitrary: it encodes \textbf{information structure}.

\begin{description}
  \item[New information] tends to come at the \textbf{end} of the sentence 
    (the ``rheme'' or focus position). This is why Polish often places the 
    subject after the verb, or saves the most concrete noun for last.
  \item[Given / known information] tends to come early.
  \item[Verb-first] is used for dramatic effect, as in: 
    \textit{Umarł Fiszer} — Fiszer died (the death before the name).
\end{description}

This freedom means that close translation often fails to capture 
the author's emphasis. Notes in the chapter analyses point out 
significant word-order choices where they carry meaning.

% ────────────────────────────────────────────────────────────
\section{Reflexive Verbs and \textit{się}}

The particle \textbf{się} is one of the most versatile elements in 
Polish. It functions as:

\begin{description}
  \item[Reflexive marker] The action reflects back on the subject: 
    \textit{zbliżyć się} — to bring oneself closer, to approach.
  \item[Reciprocal marker] \textit{spotkać się} — to meet (each other).
  \item[Impersonal/passive marker] \textit{mówi się} — one says, 
    it is said. Removes the agent entirely.
  \item[Part of the verb's inherent meaning] Some verbs simply 
    require \textit{się} with no separable reflexive meaning: 
    \textit{bać się} (to be afraid), \textit{cieszyć się} (to be glad).
\end{description}

% ────────────────────────────────────────────────────────────
\section{Participles}

Polish uses participial constructions extensively — more so than 
conversational English, but comparable to written French or Latin.

\begin{description}
  \item[Present active participle (\textit{imiesłów przymiotnikowy czynny})] 
    Describes an ongoing action of the noun it modifies. 
    Formed with suffix \textbf{-ący/-ąca/-ące}: 
    \textit{grupka \textbf{stojąca} przy grobie} — the group standing by the grave.
  \item[Past passive participle (\textit{imiesłów przymiotnikowy bierny})] 
    Describes a completed action done to the noun. 
    Formed with suffix \textbf{-ny/-na/-ne} or \textbf{-ty/-ta/-te}: 
    \textit{świeżo \textbf{wykopany} grób} — the freshly dug grave; 
    \textit{zapomniana epidemia} — the forgotten epidemic.
  \item[Adverbial participle (\textit{imiesłów przysłówkowy})] 
    Describes a simultaneous action of the subject. Does not decline. 
    Formed with suffix \textbf{-ąc}: 
    \textit{\textbf{odkręcając} nakrętkę} — while unscrewing the cap.
\end{description}

% ============================================================

\chapter*{Polish Grammar Foundations}
\addcontentsline{toc}{chapter}{Polish Grammar Foundations}
\markboth{Polish Grammar Foundations}{}

This chapter provides a compact reference for the grammatical structures 
encountered throughout the reading. It is not intended as a complete grammar 
— rather, as a touchstone for constructs that recur in the text. Cross-references 
to passages where each construct first appears are noted where relevant.

% ────────────────────────────────────────────────────────────
\section{Nouns and Gender}

Polish nouns have three grammatical genders: \textbf{masculine}, \textbf{feminine}, 
and \textbf{neuter}. Unlike French, where gender must largely be memorised, Polish 
gender is often inferable from the noun's ending.

\begin{description}
  \item[Masculine] Typically end in a consonant: \polword{grób}{grave}, 
    \polword{pogrzeb}{funeral}, \polword{syn}{son}.
  \item[Feminine] Typically end in \textbf{-a} or \textbf{-i/y}: 
    \polword{wdowa}{widow}, \polword{grupka}{small group}, 
    \polword{nadzieja}{hope}.
  \item[Neuter] Typically end in \textbf{-o} or \textbf{-e}: 
    \polword{dzieciństwo}{childhood}, \polword{dziecko}{child}.
\end{description}

\grambox{Note on animacy}{Masculine nouns are further divided into 
\textit{animate} (persons and animals) and \textit{inanimate}. This distinction 
affects the accusative case ending and is particularly important for direct objects.}

Nouns also form diminutives extensively in Polish, often with suffixes 
\textbf{-ka}, \textbf{-ko}, \textbf{-ek}, \textbf{-eczka}. These indicate 
smallness but frequently carry emotional colouring — affection, sympathy, or 
irony — rather than literal size. Examples encountered in Chapter 1: 
\polword{grupka}{little group} (from \textit{grupa}), 
\polword{garstka}{small handful} (from \textit{garść}), 
\polword{flaszka}{bottle, colloquial} (from \textit{flasza}),
\polword{podstawówka}{primary school, affectionate} (from \textit{szkoła podstawowa}).

% ────────────────────────────────────────────────────────────
\section{The Seven Cases}

Polish has seven grammatical cases. Each case encodes a grammatical 
relationship that English expresses through word order or prepositions. 
The Latin equivalents are noted as a reference point.

\begin{description}

  \item[Nominative (\textit{mianownik})] The subject of the sentence. 
    Who or what performs the action. Equivalent to Latin nominative.\\
    \textit{Groszek \textbf{powiedział}\ldots} — Groszek said\ldots

  \item[Genitive (\textit{dopełniacz})] Possession, absence, quantity, 
    and the object of many prepositions (\textit{do, od, z, bez, dla, oprócz}). 
    The most frequently used case after the nominative. Equivalent to Latin genitive.\\
    \textit{dla \textbf{uczczenia} jego \textbf{pamięci}} — for the honouring of his memory\\
    \textit{garstki \textbf{obecnych}} — a handful of those present

  \item[Dative (\textit{celownik})] The indirect object — to whom or for whom. 
    Equivalent to Latin dative. Encodes dedication, as seen in the book's 
    epigraph: \textit{\textbf{Ojcu}, \textbf{Synowi}} — to [my] father, to [my] son.

  \item[Accusative (\textit{biernik})] The direct object of most verbs, 
    and the object of motion prepositions (\textit{na, w, przez, za}).
    Equivalent to Latin accusative.\\
    \textit{Zostaw \textbf{tę flaszkę}!} — Leave that bottle!

  \item[Instrumental (\textit{narzędnik})] Means, accompaniment, and 
    identity (used after the verb \textit{być} — to be, for professions 
    and states). Also follows the preposition \textit{z} (with) and \textit{za} 
    (behind/after in positional sense). Equivalent to Latin ablative of means.\\
    \textit{z \textbf{psem}} — with a dog\\
    \textit{z \textbf{wykształcenia} inżynier} — an engineer by training

  \item[Locative (\textit{miejscownik})] Location and topic. 
    \textbf{Always} used with a preposition (\textit{w, na, przy, o, po}).
    Never appears without one. Equivalent to Latin ablative of place.\\
    \textit{na \textbf{cmentarzu}} — at the cemetery\\
    \textit{przy \textbf{grobie}} — by the grave

  \item[Vocative (\textit{wołacz})] Direct address. Used when calling 
    or speaking to someone by name or title. Relatively rare in prose 
    but common in dialogue.\\
    \textit{Mamo!} — Mum! (addressing one's mother)

\end{description}

\grambox{Key principle}{Case endings change for both nouns \textit{and} 
adjectives — they must always agree in gender, number, and case. 
This is identical to Latin adjective-noun concord: 
\textit{wczesną jesienią} (in early autumn) — both words take the 
instrumental feminine singular ending.}

% ────────────────────────────────────────────────────────────
\section{Adjectives}

Adjectives agree with their nouns in gender, number, and case. The 
ending of the adjective therefore changes to match the noun it modifies, 
across all seven cases and three genders — 42 potential forms in principle, 
though many are identical in practice.

\textit{świeżo wykopany grób} (freshly dug grave, masc. nom.) — the participle 
\textit{wykopany} agrees with \textit{grób}.

\textit{świeżo wykopanym grobie} (by the freshly dug grave, masc. loc.) — both 
adjective and noun shift to the locative.

% ────────────────────────────────────────────────────────────
\section{Verbs: Aspect}

The single most important concept in Polish verbal grammar, and the one 
with no direct Latin equivalent, is \textbf{verbal aspect}.

Every Polish verb exists in two versions:

\begin{description}
  \item[Imperfective (\textit{niedokonany})] Describes ongoing, repeated, 
    or incomplete actions. What was happening, what happens habitually, 
    what is in progress.
  \item[Perfective (\textit{dokonany})] Describes completed, bounded, 
    or one-time actions. What happened (and is done), what will definitively happen.
\end{description}

\grambox{Key contrast}{\polword{pisać}{to write (imperfective, ongoing)} 
vs.\ \polword{napisać}{to write (perfective, completed)}. The prefix 
\textit{na-} here creates the perfective. This pattern — imperfective root, 
prefixed perfective — is one of the most common in Polish.}

Aspect affects \textbf{what tenses are possible}. Perfective verbs have 
no present tense (a completed action cannot be in progress now). Their 
``present-looking'' form is actually future: \textit{napiszę} = I will write (and finish it).

\subsection*{Common perfectivising prefixes}

These prefixes convert imperfective verbs to perfective and are extremely 
productive in Polish. Recognising them unlocks large swaths of vocabulary.

\begin{description}
  \item[\textbf{za-}] Often signals the onset or completion of an action: 
    \polword{pamiętać}{to remember} → \polword{zapomnieć}{to forget}; 
    \polword{truć}{to poison} → \polword{zatruć}{to poison (completely)}.
  \item[\textbf{po-}] Often signals completion or doing something for a while: 
    \polword{czekać}{to wait} → \polword{poczekać}{to wait (for a bit / done waiting)}.
  \item[\textbf{na-}] Completion, sufficiency: 
    \polword{pisać}{to write} → \polword{napisać}{to write (and finish)}.
  \item[\textbf{od-}] Away from, completing a departure: 
    \polword{iść}{to go} → \polword{odejść}{to go away, depart}.
  \item[\textbf{wy-}] Out, thoroughly: 
    \polword{kopać}{to dig} → \polword{wykopać}{to dig out/up}.
  \item[\textbf{z- / s-}] Completion, bringing together: 
    \polword{bliski}{close} → \polword{zbliżyć się}{to draw closer}.
\end{description}

% ────────────────────────────────────────────────────────────
\section{Verb Conjugation}

Polish verbs encode person and number in their endings, making subject 
pronouns optional (as in Latin and Russian). The pronouns are used 
for emphasis or contrast.

\subsection*{Present tense — example: \textit{czekać} (to wait, imperfective)}

\begin{tabular}{lll}
  \textbf{Person} & \textbf{Singular} & \textbf{Plural} \\
  1st & czekam (I wait) & czekamy (we wait) \\
  2nd & czekasz (you wait) & czekacie (you wait) \\
  3rd & czeka (he/she/it waits) & czekają (they wait) \\
\end{tabular}

\subsection*{Past tense — gender marking}

Uniquely among European languages, Polish past tense verbs agree with 
the gender of the subject. This gender marking is a suffix added to the 
past stem:

\begin{description}
  \item[Masculine singular] \textbf{-ł}: \textit{powiedział} (he said), 
    \textit{spotkałem} (I [masc.] met)
  \item[Feminine singular] \textbf{-ła}: \textit{czekała} (she waited), 
    \textit{musiała} (she had to)
  \item[Neuter singular] \textbf{-ło}: \textit{miało} (it was supposed to)
  \item[Virile plural (masc. personal)] \textbf{-li}: \textit{zdecydowaliśmy} 
    (we [men/mixed] decided), \textit{ustawili się} (they positioned themselves)
  \item[Non-virile plural] \textbf{-ły}: used for groups of women or non-persons
\end{description}

\grambox{Virile vs.\ non-virile}{Polish distinguishes between groups 
containing at least one man (\textit{virile}, personal masculine plural) 
and all other plurals. This distinction affects verb endings, adjective 
forms, and pronouns throughout the grammar.}

% ────────────────────────────────────────────────────────────
\section{Key Impersonal Constructions}

Polish uses several impersonal constructions that have no direct English 
equivalent but will be encountered on virtually every page.

\begin{description}
  \item[\polword{trzeba}{one must, it is necessary}] Followed by an 
    infinitive. No subject. \textit{Trzeba kontynuować} — one must continue. 
    Equivalent to French \textit{il faut}.
  \item[\polword{można}{one can, it is possible}] Followed by an 
    infinitive. \textit{Można było tam dostać} — one could get there. 
    Equivalent to French \textit{on peut}.
  \item[\polword{warto}{it is worth}] Followed by an infinitive. 
    \textit{Warto wiedzieć} — it is worth knowing.
\end{description}

% ────────────────────────────────────────────────────────────
\section{Essential Particles and Connectives}

A small set of particles carries enormous communicative weight in Polish 
and will be encountered constantly.

\begin{description}
  \item[\polword{już}{already, now, anymore}] Marks a transition or 
    change of state. \textit{Już rozumiem} — I understand now (I've crossed 
    into understanding). With negation: \textit{już nie} — not anymore.
  \item[\polword{jeszcze}{still, yet, even more}] Marks continuation 
    or addition. \textit{Jeszcze jeden} — yet one more. With negation: 
    \textit{jeszcze nie} — not yet.
  \item[\polword{już nie}{not anymore}] Something was happening but has stopped.
  \item[\polword{jeszcze nie}{not yet}] Something has not happened but may.
  \item[\polword{też / także}{also, too}] \textit{Też} is conversational; 
    \textit{także} is slightly more formal.
  \item[\polword{wcale nie}{not at all}] Emphatic negation. 
    \textit{Wcale nie odchodzi} — it does not go away at all.
  \item[\polword{właśnie}{exactly, just now, precisely}] 
    Confirms or emphasises: \textit{właśnie} as a standalone response 
    means ``exactly'' or ``quite so.''
  \item[\polword{chyba}{probably, I suppose, surely}] Expresses 
    uncertainty or mild assertion. Very common in speech.
  \item[\polword{przecież}{after all, but surely}] Implies that 
    something should be obvious. \textit{Przecież wiesz} — you know 
    this perfectly well (after all).
\end{description}

% ────────────────────────────────────────────────────────────
\section{Conjunctions: Because}

Three words translate as ``because'' with different registers:

\begin{description}
  \item[\polword{bo}{because}] Casual, spoken, everyday. 
    The default in conversation.
  \item[\polword{ponieważ}{because}] Standard written and formal spoken. 
    Neutral and widely applicable.
  \item[\polword{gdyż}{because/for}] Literary and elevated. 
    Signals deliberate, crafted prose. Equivalent to French \textit{car}.
\end{description}

% ────────────────────────────────────────────────────────────
\section{Word Order}

Polish word order is considerably freer than English, because case 
endings — not position — identify grammatical role. However, order 
is far from arbitrary: it encodes \textbf{information structure}.

\begin{description}
  \item[New information] tends to come at the \textbf{end} of the sentence 
    (the ``rheme'' or focus position). This is why Polish often places the 
    subject after the verb, or saves the most concrete noun for last.
  \item[Given / known information] tends to come early.
  \item[Verb-first] is used for dramatic effect, as in: 
    \textit{Umarł Fiszer} — Fiszer died (the death before the name).
\end{description}

This freedom means that close translation often fails to capture 
the author's emphasis. Notes in the chapter analyses point out 
significant word-order choices where they carry meaning.

% ────────────────────────────────────────────────────────────
\section{Reflexive Verbs and \textit{się}}

The particle \textbf{się} is one of the most versatile elements in 
Polish. It functions as:

\begin{description}
  \item[Reflexive marker] The action reflects back on the subject: 
    \textit{zbliżyć się} — to bring oneself closer, to approach.
  \item[Reciprocal marker] \textit{spotkać się} — to meet (each other).
  \item[Impersonal/passive marker] \textit{mówi się} — one says, 
    it is said. Removes the agent entirely.
  \item[Part of the verb's inherent meaning] Some verbs simply 
    require \textit{się} with no separable reflexive meaning: 
    \textit{bać się} (to be afraid), \textit{cieszyć się} (to be glad).
\end{description}

% ────────────────────────────────────────────────────────────
\section{Participles}

Polish uses participial constructions extensively — more so than 
conversational English, but comparable to written French or Latin.

\begin{description}
  \item[Present active participle (\textit{imiesłów przymiotnikowy czynny})] 
    Describes an ongoing action of the noun it modifies. 
    Formed with suffix \textbf{-ący/-ąca/-ące}: 
    \textit{grupka \textbf{stojąca} przy grobie} — the group standing by the grave.
  \item[Past passive participle (\textit{imiesłów przymiotnikowy bierny})] 
    Describes a completed action done to the noun. 
    Formed with suffix \textbf{-ny/-na/-ne} or \textbf{-ty/-ta/-te}: 
    \textit{świeżo \textbf{wykopany} grób} — the freshly dug grave; 
    \textit{zapomniana epidemia} — the forgotten epidemic.
  \item[Adverbial participle (\textit{imiesłów przysłówkowy})] 
    Describes a simultaneous action of the subject. Does not decline. 
    Formed with suffix \textbf{-ąc}: 
    \textit{\textbf{odkręcając} nakrętkę} — while unscrewing the cap.
\end{description}


% ── Book Chapters ────────────────────────────────────────────
% Karta tytułowa (front matter passages)
% ============================================================
%  ch00_karta_tytulowa.tex  —  Front Matter Passages
%  Dedication, epigraphs, and title-page text
% ============================================================

\chapter*{Karta tytułowa}
\addcontentsline{toc}{chapter}{Karta tytułowa}
\markboth{Karta tytułowa}{}

% ── Opening line ─────────────────────────────────────────────
\section{Motto}

\poltext{Zatrute dzieciństwo nie pozwoli zapomnieć.}
\poltrans{A poisoned childhood will not allow one to forget.}

\textbf{Zatrute} — \textit{poisoned} (adjective from \textit{zatruć}, to poison). 
Neuter nominative singular, agreeing with \textit{dzieciństwo}. 
Adjectival agreement exactly as in Latin.

\textbf{dzieciństwo} — \textit{childhood} (neuter noun, nominative singular, subject). 
The suffix \textit{-stwo} is a productive Polish nominaliser, like English 
\textit{-hood} or \textit{-ness}. Root: \textit{dziecko} (child).

\textbf{nie pozwoli} — \textit{will not allow}. Third person singular future 
perfective of \textit{pozwolić}. The perfective aspect signals a definitive, 
completed future action — not an ongoing permission withheld, but an absolute refusal.

\textbf{zapomnieć} — \textit{to forget} (perfective infinitive). 
From \textit{zapominać/zapomnieć}. The \textit{za-} prefix perfectivises 
\textit{pamiętać} (to remember) and reverses its meaning: to have-remembered-away.

\grambox{Impersonal construction}{\textit{nie pozwoli zapomnieć} uses an 
impersonal infinitive — there is no explicit object (one, you, we). This 
gives the line a universal quality. The book's subtitle is 
\textit{Zapomniana epidemia} (the forgotten epidemic); this opening line 
directly inverts that word: the epidemic may be forgotten, but a poisoned 
childhood will not \textit{permit} forgetting.}

% ── Dedication ───────────────────────────────────────────────
\section{Dedication}

\poltext{Ojcu, Hubertowi\\Synowi, Michałowi}
\poltrans{To [my] father, Hubert\\To [my] son, Michał}

Both lines use the \textbf{dative case} — the case of giving and dedicating. 
English requires the preposition ``to''; Polish encodes it in the noun ending.

\polword{Ojciec}{father} → \polword{Ojcu}{to [my] father} (dative singular).

\polword{Syn}{son} → \polword{Synowi}{to [my] son} (dative singular).

Proper names follow exactly the same pattern: 
\polword{Hubert}{} → \polword{Hubertowi}{}; 
\polword{Michał}{} → \polword{Michałowi}{}. 
Even personal names decline fully in Polish — a feature that would feel 
strange in English but is entirely natural here, and is identical to Latin.

The parallel structure — father and son, two generations — mirrors each 
other perfectly in form as well as meaning. No verb, no elaboration. 
The author dedicates a book about the damage done to children to his own 
father and his own son. That juxtaposition is quietly devastating.

% ── Stephen King epigraph ────────────────────────────────────
\section{Epigraph}

\poltext{Czasami to, o czym chcielibyśmy zapomnieć, co miało odejść 
w przeszłość, wcale nie odchodzi.}
\poltrans{Sometimes that which we would like to forget, that which was 
supposed to recede into the past, does not go away at all.}

\textbf{Czasami} — \textit{sometimes}. A gentle, resigned opening.

\textbf{to, o czym\ldots} — \textit{that which\ldots} Relative clause construction. 
\textit{To} is the antecedent; \textit{o czym} is the relative pronoun in the 
locative (literally ``about which'').

\textbf{chcielibyśmy zapomnieć} — \textit{we would like to forget}. 
\textit{Chcielibyśmy} is the conditional first person plural of \textit{chcieć} 
(to want). The \textit{-by-} infix is the Polish conditional marker, equivalent 
to French \textit{-rais/-rait}.

\textbf{co miało odejść} — \textit{that which was supposed to go away}. 
\textit{Miało} is the neuter past of \textit{mieć} used as an auxiliary 
meaning ``was supposed to / was meant to.'' \textit{Odejść}: the prefix 
\textit{od-} (away from) combined with \textit{iść} (to go).

\textbf{w przeszłość} — \textit{into the past}. \textit{Przeszłość} in the 
accusative, required by the motion of \textit{odejść} — going \textit{into} 
somewhere. Root: \textit{przeszły} (past, gone).

\textbf{wcale nie} — \textit{not at all}. \textit{Wcale} is an intensifier 
of negation. Compare \textit{nie odchodzi} (doesn't go away) with 
\textit{wcale nie odchodzi} (doesn't go away at all — simply refuses to leave).

\textbf{odchodzi} — \textit{goes away} (imperfective present). The imperfective 
aspect emphasises the ongoing, continuous failure to leave.

The sentence is structured as a long double relative clause — two layers of 
``that which we wanted gone'' — before the short, flat conclusion: 
\textit{wcale nie odchodzi}. Note the echo of \textit{zapomnieć} from the 
very first line of the book.

\vspace{1em}

\poltext{Czasami to do ciebie wraca.\\{\small --- Stephen King, \textit{To}}}
\poltrans{Sometimes it comes back to you.\\--- Stephen King, \textit{It}}

\textbf{to} — \textit{it} — doing double duty. A neutral demonstrative pronoun 
in Polish, but also the Polish title of King's novel \textit{It}. The sentence 
is simultaneously a quote about unresolvable horror and a quiet nod to the 
book it comes from.

\textbf{do ciebie} — \textit{to you}. \textit{Do} takes the genitive; 
\textit{ciebie} is the emphatic genitive/accusative of \textit{ty} (you). 
The emphatic form rather than the shorter \textit{do cię} gives it 
deliberate, personal weight.

\textbf{wraca} — \textit{returns, comes back}. Imperfective present, third 
person singular. The imperfective aspect: it doesn't just return once, 
it keeps returning.

The author opens with his own sentence about things that refuse to leave, 
then follows with King — whose novel \textit{It} is literally about childhood 
trauma that was supposed to stay buried but keeps returning. The Polish title 
\textit{To} makes the pronoun in the sentence indistinguishable from the 
title itself. After the previous sentence's elegiac tone, this lands with 
quiet menace.

\grambox{Już vs.\ jeszcze — first encounter}{\textbf{Już} marks a transition 
or arrival: \textit{już jest} — he's here now; \textit{już nie} — not anymore. 
\textbf{Jeszcze} marks continuation or addition: \textit{jeszcze jest} — he's 
still here; \textit{jeszcze nie} — not yet. Latin parallel: \textit{już} 
resembles \textit{iam}; \textit{jeszcze} resembles \textit{adhuc}.}


% Chapter 1
% ============================================================
%  ch01_jeszcze_jeden_pogrzeb.tex  —  Chapter 1
%  Jeszcze jeden pogrzeb (pp. 5–30)
% ============================================================

\chapter{Jeszcze jeden pogrzeb}
\markboth{Jeszcze jeden pogrzeb}{}

% ── Chapter title ────────────────────────────────────────────
\section{Tytuł rozdziału / Chapter Title}

\poltext{Jeszcze jeden pogrzeb}
\poltrans{Yet Another Funeral / One More Funeral}

\textbf{Jeszcze} — \textit{still, yet, another}. Implies continuation and 
accumulation, even weariness. \textit{Jeszcze raz} — once more; 
\textit{jeszcze jeden} — yet one more, with the implication that there 
have already been others.

\textbf{jeden} — \textit{one} (masculine nominative, agreeing with \textit{pogrzeb}).

\textbf{pogrzeb} — \textit{funeral} (masculine noun, nominative). 
Root: \textit{grzebać} — to bury, to dig. The earthiness of the root word 
is audible in it.

The devastating efficiency of this title is in \textit{jeszcze}. Not ``a funeral'' 
or even ``another funeral'' — but \textit{yet one more}, the way someone says 
it who has lost count.

% ── First paragraph ──────────────────────────────────────────
\section{Pierwsze akapity / Opening Paragraphs}

\poltext{Umarł Fiszer, kolega z podstawówki. Postać ważna, gdyż jego dziadek, 
stary Fiszer, miał kiosk. Można było tam dostać spod lady komiksy z serii: 
Kapitan Kloss, Kapitan Żbik albo Podziemny Front --- dla nas absolutne rarytasy, 
a także czerwone Carmeny i niebieskie Caro, towar jeszcze bardziej pożądany, 
zważywszy, że były to pierwsze w życiu papierosy palone przez nas w krzakach. 
Już w piątej klasie.}
\poltrans{Fiszer died --- a classmate from primary school. An important figure, 
because his grandfather, old Fiszer, had a kiosk. You could get there under the 
counter comics from the series: Kapitan Kloss, Kapitan Żbik, or Podziemny Front 
--- for us absolute rarities, and also red Carmens and blue Caros, goods even more 
desirable considering that these were the first cigarettes smoked by us in the 
bushes. Already in fifth grade.}

\textbf{Umarł Fiszer} — verb first, subject second. Polish allows this inversion 
freely; the death lands before we know who died.

\textbf{kolega z podstawówki} — \textit{podstawówka} is the colloquial, affectionate 
diminutive of \textit{szkoła podstawowa} (primary school). The suffix \textit{-ka} 
sets the nostalgic, intimate tone immediately.

\textbf{gdyż} — \textit{because} (literary register). Signals deliberate craft. 
See Grammar Foundations §\,Conjunctions for the full \textit{bo / ponieważ / gdyż} 
register spectrum.

\textbf{Można było tam dostać} — \textit{one could get there}. Impersonal 
\textit{można} (one can) + \textit{było} (past auxiliary) + 
\textit{dostać} (to obtain, perfective infinitive).

\textbf{spod lady} — \textit{from under the counter}. \textit{Spod} is a 
compound preposition (from under), taking the genitive. Under-the-counter goods 
in communist-era Poland were a cultural institution.

\textbf{rarytasy} — \textit{rarities, treats}. Plural of \textit{rarytas} — 
an old-fashioned, savoured word.

\textbf{czerwone Carmeny i niebieskie Caro} — Polish cigarette brands identified 
entirely by their packaging colour. \textit{Czerwone} and \textit{niebieskie} are 
plural adjectives in agreement with the brand names.

\textbf{zważywszy, że} — \textit{considering that, given that}. A formal 
participial construction applied with dry irony to children smoking in bushes.

\textbf{palone przez nas w krzakach} — \textit{smoked by us in the bushes}. 
Passive participial construction. \textit{Krzaki} (bushes) — a word that 
needs no translation once heard.

\textbf{Już w piątej klasie.} — \textit{Already in fifth grade.} The fragment 
stands alone. \textit{Już} marks arrival at a milestone — fifth grade, roughly 
age eleven. The author's eyebrow is raised at his own childhood.

% ────────────────────────────────────────────────────────────

\poltext{Fiszer umarł nagle, niespodziewanie i kłopotliwie.}
\poltrans{Fiszer died suddenly, unexpectedly, and inconveniently.}

Three adverbs in escalating syllable count, each adding weight before 
the final, jarring word.

\textbf{nagle} — \textit{suddenly}. Abrupt, no warning.

\textbf{niespodziewanie} — \textit{unexpectedly}. From \textit{spodziewać się} 
(to expect); \textit{nie-} prefix negates it, exactly as English \textit{un-expected}.

\textbf{kłopotliwie} — \textit{inconveniently, troublesomely}. From 
\textit{kłopot} (trouble, bother, hassle). Applied to a man's death, 
this is either bleakly comic or quietly brutal. The author could have 
written \textit{tragicznie} (tragically). He chose \textit{kłopotliwie}. 
That choice announces the voice of the book.

% ────────────────────────────────────────────────────────────

\poltext{Zrozpaczona wdowa musiała transportować jego zwłoki z południa Francji.}
\poltrans{The grief-stricken widow had to transport his remains from the south of France.}

\textbf{Zrozpaczona} — \textit{grief-stricken}. From \textit{rozpacz} (despair). 
The \textit{z-} prefix indicates a state fully arrived at. A strong, unambiguous word.

\textbf{musiała} — \textit{had to, was obliged to}. Past tense of \textit{musieć}. 
Not ``chose to'' or ``arranged to'' but \textit{had to} — the logistics were inescapable.

\textbf{transportować} — borrowed from international vocabulary. The clinical, 
bureaucratic register against \textit{zrozpaczona wdowa} is quietly jarring: 
grief on one side, freight paperwork on the other.

\textbf{zwłoki} — \textit{mortal remains, corpse}. Exists only in the plural 
in Polish, like English ``remains.'' More formal than \textit{trup} (body), 
but unflinching.

\textbf{z południa Francji} — \textit{from the south of France}. 
\textit{Południe} means both \textit{south} and \textit{noon} — 
sharing a root with the sun's position. Genitive follows \textit{z} (from).

The contrast with \textit{kłopotliwie} is now fully clear: the 
``inconvenience'' mentioned so dryly has a very human cost, borne entirely 
by the widow.

% ────────────────────────────────────────────────────────────

\poltext{Zdecydowaliśmy, że dla uczczenia jego pamięci trzeba się napić, 
i z takim postanowieniem udaliśmy się na pogrzeb.}
\poltrans{We decided that to honour his memory one must have a drink, and 
with that resolution we set off for the funeral.}

\textbf{Zdecydowaliśmy} — \textit{we decided}. Perfective past, first person 
plural. A decision fully made.

\textbf{dla uczczenia jego pamięci} — \textit{for the honouring of his memory}. 
Formal, almost ceremonial. \textit{Uczczenie} is a verbal noun from 
\textit{uczcić} (to honour, to commemorate). The genitive chain 
\textit{uczczenia jego pamięci} illustrates how Polish stacks nouns in the genitive.

\textbf{trzeba się napić} — \textit{one must have a drink}. 
\textit{Trzeba} (one must) reappears from the epigraph material. 
\textit{Napić się}: the reflexive \textit{się} with \textit{napić} means 
to have a drink, to drink one's fill — a distinction between an activity 
(\textit{pić}, to drink) and a completed, satisfying act (\textit{napić się}).

\textbf{z takim postanowieniem} — \textit{with such a resolution}. 
Instrumental after \textit{z} (with). \textit{Postanowienie} is a wonderfully 
weighty word for what is essentially ``let's get drunk at the funeral.''

\textbf{udaliśmy się} — \textit{we set off, we betook ourselves}. Slightly 
elevated register for going somewhere. The formality is faintly comic here.

The sentence wraps a decision to drink in the language of solemn civic duty. 
The gap between the gravity of the language and the actual resolution is where 
all the humour lives.

% ────────────────────────────────────────────────────────────

\poltext{Na cmentarzu w Dąbrówce Małej przy świeżo wykopanym grobie czekała 
na kondukt grupka żałobników.}
\poltrans{At the cemetery in Dąbrówka Mała, by the freshly dug grave, a small 
group of mourners waited for the funeral procession.}

\textbf{Na cmentarzu} — \textit{at the cemetery}. \textit{Na} with the locative 
for location. \textit{Cmentarz} — a word worth knowing immediately in Poland, 
where cemeteries are extraordinary cultural spaces.

\textbf{w Dąbrówce Małej} — \textit{in Dąbrówka Mała}. A small district of 
Katowice in Silesia. The place name declines: \textit{Dąbrówka Mała} 
(nominative) → \textit{Dąbrówce Małej} (locative after \textit{w}). 
Both words change their endings in agreement — noun and adjective declining 
together, exactly as in Latin.

\textbf{przy świeżo wykopanym grobie} — \textit{by the freshly dug grave}. 
\textit{Przy} takes the locative. \textit{Świeżo} (adverb: freshly); 
\textit{wykopanym} (past passive participle, locative: dug). Note the echo 
of \textit{Groby są świeże} from the epigraph — the author has made good 
on his own promise.

\textbf{grupka żałobników} — \textit{a small group of mourners}. 
\textit{Grupka} — diminutive of \textit{grupa}, the \textit{-ka} suffix 
quietly measuring how few people are present. \textit{Żałobników} is the 
genitive plural of \textit{żałobnik} (mourner), required after \textit{grupka}.

The sentence moves the camera slowly inward: location → specific location 
→ the grave → the waiting → the handful of people.

% ────────────────────────────────────────────────────────────

\poltext{Za nimi, przy głównej alejce, ustawili się dyskretnie grabarze.}
\poltrans{Behind them, by the main path, the gravediggers had discreetly 
taken their positions.}

\textbf{Za nimi} — \textit{behind them}. \textit{Za} with the instrumental 
for position. \textit{Nimi} is the instrumental form of \textit{oni/one} (they).

\textbf{przy głównej alejce} — \textit{by the main path}. 
\textit{Alejka} is a diminutive of \textit{aleja} — a little lane between graves.

\textbf{ustawili się} — \textit{positioned themselves}. 
\textit{Ustawić się}: reflexive perfective — they have arranged themselves 
and are settled.

\textbf{dyskretnie} — \textit{discreetly}. From French/Latin \textit{discret}. 
Its position between verb and subject gives it quiet emphasis.

\textbf{grabarze} — \textit{gravediggers}. From \textit{grób} (grave) — 
the same root recurring throughout the chapter. Placed last in the sentence: 
the most concrete image saved for the end.

The discreet gravediggers waiting in the wings carry an unspoken 
practicality that echoes the earlier \textit{kłopotliwie}. Everyone 
is performing their assigned role.

% ────────────────────────────────────────────────────────────

\poltext{Wczesną jesienią, ciepłą i kolorową, pasującą bardziej do spaceru 
z psem lub grzybobrania niż do śmierci, na cmentarzu parafialnym kościoła 
pod wezwaniem Świętego Antoniego, patrona między innymi od spraw beznadziejnych, 
po prawie czterdziestu latach spotkałem szkolne koleżanki i kolegów z podstawówki.}
\poltrans{In the early autumn, warm and colourful, more suited to a walk with 
a dog or mushroom picking than to death, at the parish cemetery of the church 
under the patronage of Saint Anthony, patron among other things of hopeless 
causes, after nearly forty years I met my schoolgirls and schoolboys from 
primary school.}

\textbf{Wczesną jesienią} — \textit{in early autumn}. Instrumental case used 
for time periods. Both words decline together.

\textbf{pasującą bardziej do\ldots niż do} — \textit{more suited to\ldots than to}. 
\textit{Bardziej\ldots niż} is the standard Polish comparative construction. 
\textit{Pasować do} (to suit, to fit) always takes the genitive.

\textbf{grzybobrania} — \textit{mushroom picking}. A compound: 
\textit{grzyb} (mushroom) + \textit{branie} (taking/picking). 
Mushroom foraging is a deeply serious Polish cultural institution, and its 
appearance here as the opposite of death is both funny and characteristically Polish.

\textbf{pod wezwaniem Świętego Antoniego} — \textit{under the patronage of 
Saint Anthony}. Literally ``under the invocation of.'' The fixed Polish phrase 
for a church's dedicatory patron.

\textbf{patrona między innymi od spraw beznadziejnych} — \textit{patron among 
other things of hopeless causes}. Saint Anthony is the patron of lost things 
and hopeless cases. \textit{Między innymi} (among other things) is used 
constantly in Polish. \textit{Beznadziejnych}: from \textit{nadzieja} (hope), 
negated by \textit{bez-} (without).

\textbf{po prawie czterdziestu latach} — \textit{after nearly forty years}. 
\textit{Po} with the locative for ``after.'' \textit{Czterdziestu} is the 
genitive plural of \textit{czterdzieści} (forty), required after \textit{prawie} (almost).

\textbf{spotkałem szkolne koleżanki i kolegów} — \textit{I met schoolgirls 
and schoolboys}. First appearance of a first-person singular verb in the 
narrative. \textit{Spotkałem}: perfective past, masculine singular. Note that 
Polish distinguishes \textit{koleżanka} (female classmate) and \textit{kolega} 
(male classmate) — both are needed to include everyone.

The sentence delays the main verb \textit{spotkałem} through layer after layer 
of atmospheric detail. The weather is wrong for a funeral. The church's own 
patron saint specialises in hopeless causes. By the time we arrive at the actual 
event, it has been framed as something faintly absurd, faintly doomed, and entirely human.

% ────────────────────────────────────────────────────────────

\poltext{Oprócz garstki obecnych wspomnijmy też całą resztę: nieobecnych lub martwych.}
\poltrans{Besides the handful of those present, let us also remember all the rest: 
the absent or the dead.}

\poltext{Najpierw my, obecni.}
\poltrans{First, us --- the ones who are here.}

\textbf{Oprócz} — \textit{besides, apart from}. Takes the genitive.

\textbf{garstki} — genitive singular of \textit{garstka} — a handful, literally 
the amount that fits in a cupped hand. A second diminutive quietly measuring 
how few remain.

\textbf{wspomnijmy} — \textit{let us remember}. First person plural imperative 
of \textit{wspomnieć}. The \textit{-my} ending is the inclusive ``let us'' form. 
The author pulls the reader into the obligation.

\textbf{całą resztę} — \textit{all the rest}. Accusative after \textit{wspomnijmy}.

\textbf{nieobecnych lub martwych} — \textit{the absent or the dead}. Genitive 
plural in apposition to \textit{resztę}. The colon before this phrase makes it 
land as a definition: the rest = the absent or the dead. Only two categories remain.

\textbf{Najpierw my, obecni.} — \textit{First, us --- the present ones.} 
No verb. \textit{My} (we/us) stated explicitly is emphatic in Polish, where 
the subject is normally encoded in the verb ending alone. Saying \textit{my} 
means: \textit{us, specifically, as distinct from them.}

% ────────────────────────────────────────────────────────────

\poltext{Zbliżyliśmy się do pozostałej grupki stojącej przy grobie.}
\poltrans{We drew closer to the remaining little group standing by the grave.}

\textbf{Zbliżyliśmy się} — \textit{we drew closer}. From \textit{zbliżyć się}, 
the perfective reflexive of \textit{bliski} (close). A gentle, hesitant motion 
— not \textit{podeszliśmy} (we walked up to) but a gradual closing of distance.

\textbf{do pozostałej grupki} — \textit{to the remaining little group}. 
\textit{Pozostałej}: genitive of \textit{pozostały} (remaining, left over) 
— a quietly melancholy adjective. \textit{Grupki}: the same diminutive 
as before, insisting again on their fewness.

\textbf{stojącej przy grobie} — \textit{standing by the grave}. 
Present active participle, genitive singular agreeing with \textit{grupki}.

After the rhetorical force of \textit{Najpierw my, obecni}, the narrative 
simply resumes. The mundane physical action after the weightier reflection 
is characteristic of this author's rhythm.

% ── Groszek's entrance ───────────────────────────────────────
\section{Groszek}

\poltext{--- Coś nie za dużo nas --- powiedział Groszek, odkręcając nakrętkę.}
\poltrans{`Not too many of us, is there,' said Groszek, unscrewing the cap.}

\textbf{Coś nie za dużo nas} — \textit{colloquial and understated}. \textit{Coś} 
here is not ``something'' but a softening particle — a shrug built into the grammar. 
\textit{Nie za dużo}: not too many. The whole phrase is casual, almost throwaway, 
for an observation about how many of their childhood friends are still alive.

\textbf{Groszek} — literally \textit{little pea} or \textit{little coin} 
(diminutive of \textit{grosz}, a small coin; or \textit{groch}, pea). 
A childhood nickname that has clearly followed someone into late middle age.

\textbf{odkręcając nakrętkę} — \textit{unscrewing the cap}. 
Present adverbial participle: simultaneous action. \textit{Odkręcać} 
(to unscrew, to twist open) and \textit{nakrętka} (cap, lid, screw-top) 
share the root \textit{kręcić} (to turn, to twist). He is already opening 
the bottle while speaking.

% ────────────────────────────────────────────────────────────

\poltext{--- Zostaw tę flaszkę!}
\poltrans{`Put that bottle down!'}

\textbf{Zostaw} — \textit{leave it, put it down}. Imperative singular of 
\textit{zostawić}. Sharp, direct, unambiguous.

\textbf{tę} — accusative singular feminine of the demonstrative \textit{ta} 
(this/that). Polish demonstratives decline fully; \textit{tę} is a form 
worth recognising on sight.

\textbf{flaszkę} — accusative singular of \textit{flaszka}. A colloquial, 
slightly affectionate word for a bottle — more informal than \textit{butelka}. 
The diminutive flavour makes the outrage slightly comic.

% ────────────────────────────────────────────────────────────

\poltext{Groszek, z wykształcenia inżynier, z przekonania śląski autonomista 
i w mediach społecznościowych radykalny zwolennik teorii o wyższości silników 
diesla nad elektrycznymi, posłuchał chóralnej sugestii. Na razie.}
\poltrans{Groszek, an engineer by training, a Silesian autonomist by conviction, 
and on social media a radical proponent of the theory of the superiority of 
diesel engines over electric ones, heeded the collective suggestion. For now.}

\textbf{z wykształcenia\ldots z przekonania} — \textit{by education\ldots by conviction}. 
A fixed Polish construction for introducing someone's professional identity 
and worldview simultaneously. Clean and elegant.

\textbf{śląski autonomista} — \textit{Silesian autonomist}. Silesian regionalism 
and greater autonomy from Warsaw is a genuine, lively political position. 
This detail tells you immediately who Groszek is and where he comes from.

\textbf{w mediach społecznościowych} — \textit{on social media}. 
\textit{Media społecznościowe} is the standard Polish term, a calque of 
``social media.'' Locative plural after \textit{w}.

\textbf{radykalny zwolennik teorii o wyższości\ldots nad\ldots} — 
\textit{a radical proponent of the theory of the superiority of\ldots over\ldots} 
\textit{O wyższości}: from \textit{wysoki} (high/superior); \textit{nad} 
with the instrumental expresses ``over, above.'' The specificity of the 
diesel-over-electric conviction is perfect character comedy: not just an 
opinion, but a \textit{theory}, held \textit{radically}, prosecuted on social media.

\textbf{posłuchał chóralnej sugestii} — \textit{heeded the collective suggestion}. 
\textit{Posłuchać} (to listen to, to obey) takes the genitive. 
\textit{Chóralnej}: \textit{choral, unanimous, in chorus} — from \textit{chór} 
(choir). Everyone told him to put the bottle down at once; the word chosen 
is musical, theatrical.

\textbf{Na razie.} — \textit{For now.} Two words doing enormous work. 
He complied — but only provisionally. \textit{Na razie} is one of the most 
useful phrases in Polish: ``for the time being.'' Here it functions as both 
narrative aside and character revelation. Groszek has been introduced, 
assessed, and shown to be entirely himself in one sentence.

% ── Social anxiety and the bottle ───────────────────────────
\section{Chyba nie powinno się pić}

\poltext{Chyba nie powinno się pić na cmentarzu, przynajmniej w tym kraju 
i w tym obrządku. Nawet w dobrej intencji.}
\poltrans{One probably shouldn't drink at a cemetery, at least in this country 
and in this rite. Even with good intentions.}

\textbf{Chyba} --- \textit{probably, I suppose}. Softens the observation into 
a reasonable social reflection rather than a scolding. See Grammar Foundations.

\textbf{nie powinno się pić} --- \textit{one shouldn't drink}. \textbf{Powinno} 
is the neuter form of \textit{powinien} (should) --- neuter because there is no 
named subject. \textbf{Się} makes it impersonal: the mild reproach is distributed 
across everyone present, including the speaker.

\textbf{przynajmniej} --- \textit{at least}. High-frequency hedging word.

\textbf{obrządek} --- \textit{rite, religious observance} --- specifically 
religious and formal. The speaker acknowledges the rule might differ elsewhere: 
other countries, other funerary customs. A notably cosmopolitan hedge.

\textbf{Nawet w dobrej intencji.} --- \textit{Even with good intentions.} 
A verbless fragment. \textbf{Nawet} (\textit{even}) is a high-frequency particle. 
Even granting honourable intent, it still probably shouldn't happen. Probably.

% ── Warsaw ──────────────────────────────────────────────────
\section{Słoik z Warszawy}

\poltext{Poza tym, co powiedzą nasi dawno niewidziani znajomi, jak nas 
zobaczą\ldots? Pewno nic. Zestarzeliśmy się. To będzie największą atrakcją.}
\poltrans{Besides, what will our long-unseen acquaintances say when they 
see us\ldots? Probably nothing. We've grown old. That will be the greatest attraction.}

\textbf{Poza tym} --- \textit{besides, apart from that} --- conversational 
connector, slightly more colloquial than \textit{oprócz}.

\textbf{dawno niewidziani znajomi} --- \textit{long-unseen acquaintances}. 
\textbf{Dawno niewidziani}: \textit{dawno} (long ago) modifying the past 
passive participle \textit{niewidziany}. \textbf{Znajomi} --- warmer than 
\textit{obcy} (strangers), cooler than \textit{przyjaciele} (friends).

\textbf{Pewno nic.} --- \textit{Probably nothing.} Colloquial variant of 
\textit{pewnie}. One devastating word in answer to the question just posed.

\textbf{Zestarzeliśmy się.} --- \textit{We've grown old.} The prefix \textbf{ze-} 
perfectivises \textit{starzeć się}. Stated flatly, without self-pity.

\textbf{największą atrakcją} --- superlative instrumental, as predicates after 
\textit{być} take the instrumental. The spectacle of ageing classmates will 
be the main event --- not the funeral itself.

% ────────────────────────────────────────────────────────────

\poltext{Jako „słoik" przybyły z Warszawy byłem obiektem szczególnego zainteresowania.}
\poltrans{As a `jar' who had come from Warsaw, I was an object of particular interest.}

\textbf{słoik} --- Warsaw slang for someone who has moved to the capital from 
the provinces, stereotypically arriving each Monday with a jar of home-cooked 
food. Applied to himself with ironic self-awareness: he has returned to Silesia 
as the outsider.

\textbf{obiektem szczególnego zainteresowania} --- \textit{an object of particular 
interest}: instrumental after \textit{być}, genitive chain following.

% ────────────────────────────────────────────────────────────

\poltext{Nic tak nie pobudza do, nazwijmy to, wyrażania zdania odrębnego jak Warszawa.}
\poltrans{Nothing provokes, let us call it, the expression of a dissenting 
opinion quite like Warsaw.}

\textbf{nazwijmy to} --- \textit{let us call it} --- first person plural imperative, 
a parenthetical flagging that what follows is a deliberate euphemism.

\textbf{wyrażania zdania odrębnego} --- \textit{the expression of a dissenting opinion}. 
\textbf{Odrębny} --- \textit{separate, distinct, dissenting} --- formal, almost 
legalistic. The whole phrase wraps strong regional resentment in mock-bureaucratic diction.

% ────────────────────────────────────────────────────────────

\poltext{Łączy całą Polskę w wyrażaniu jednolitego, solidarnego, wspólnego 
ponad podziałami uczucia nienawiści.}
\poltrans{It unites all of Poland in the expression of a uniform, solidary, 
shared-across-divisions feeling of hatred.}

\textbf{ponad podziałami} --- \textit{above/across divisions} --- lifted 
directly from Polish political reconciliation rhetoric, deployed with lethal 
irony. The sentence is built in the exact syntactic shape of a patriotic 
proclamation and delivers \textit{hatred} as the payload.

% ── The ceremony approaches ──────────────────────────────────
\section{Zaraz się zacznie}

\poltext{Zaraz się zacznie. Nadchodzi moja młodsza siostra z koleżanką i kolegą.}
\poltrans{It's about to start. My younger sister is approaching with a female 
friend and a male friend.}

\textbf{Zaraz} --- one of the essential Polish time words: \textit{in a moment, 
right away, any second now}. After three sentences of sardonic Warsaw commentary, 
the narrative snaps back to the immediate present.

\textbf{Nadchodzi} --- imperfective present of \textit{nadchodzić}: the prefix 
\textit{nad-} carries a sense of arrival from a distance.

% ────────────────────────────────────────────────────────────

\poltext{Wcześniej zdążyłem dowiedzieć się, że koleżanka jest świeżo po 
rozwodzie, a kolegę zobaczę o kulach.}
\poltrans{Earlier I had managed to find out that the female friend is freshly 
divorced, and that I would see the male friend on crutches.}

\textbf{zdążyłem} --- perfective past of \textit{zdążyć}: doing something 
just in time, managing to accomplish something within a window of opportunity.

\textbf{świeżo po rozwodzie} --- \textit{freshly divorced}. \textbf{Świeżo}: 
the same root as \textit{świeżo wykopanym grobie} from the opening --- 
applied here to a divorce rather than a grave, with no less bluntness.

\textbf{o kulach} --- \textit{on crutches} --- \textit{o} with the locative 
in this fixed expression.

% ── Dialogue: greetings ──────────────────────────────────────
\section{Powitania / Greetings}

\poltext{--- Cześć, siostra --- przywitałem się z jedyną z tej grupy osobą, 
z którą miałem kontakt przez te lata.}
\poltrans{`Hi, sis,' I greeted the only person from this group with whom 
I had been in contact over these years.}

\textbf{przez te lata} --- \textit{over these years} --- \textit{przez} with 
the accusative for duration. The narration quietly reveals that of everyone 
here, his sister is the only one he stayed in contact with.

% ────────────────────────────────────────────────────────────

\poltext{--- Cześć, brat.}
\poltrans{`Hi, bro.'}

She mirrors his greeting exactly. Two words --- and the warmth lands after 
everything that preceded: the funeral, old age, crutches, divorce, Warsaw 
resentment, a bottle at a graveside. The author knows when to stop.

% ────────────────────────────────────────────────────────────

\poltext{--- Cześć, Daisy. Cześć, Hanka --- Groszek też się przywitał.}
\poltrans{`Hi, Daisy. Hi, Hanka,' --- Groszek greeted them too.}

\textbf{Daisy} sits unglossed beside \textit{Hanka} --- a standard Polish 
diminutive --- without comment. Whether by birth, adoption, or some story 
not yet told, someone in this childhood group from the Katowice suburbs 
is called Daisy. Groszek, meanwhile, is performing social normalcy. 
\textit{Na razie} is still in effect.

% ── Middle age ───────────────────────────────────────────────
\section{Mielizna średniego wieku / The Sandbank of Middle Age}

\poltext{Jeżeli najważniejsze jest pierwsze wrażenie, a --- pocieszmy się --- 
nie jest, to powinniśmy być w szoku, oceniając nasz aktualny wygląd.}
\poltrans{If first impressions are what matter most --- and, let us console 
ourselves, they are not --- then we should be in shock, assessing our current appearance.}

\textbf{pocieszmy się} --- \textit{let us console ourselves} --- first person 
plural imperative, a parenthetical dropped into the conditional. The consolation 
is bleak: first impressions don't matter. Small mercies.

\textbf{aktualny} --- a Polish false friend: does not mean \textit{actual} 
in the English sense but \textit{current, as things stand now}.

% ────────────────────────────────────────────────────────────

\poltext{Czas zatarł podział na trochę starszych i trochę młodszych, wyrzucając 
nas wszystkich na tak kiedyś odległą mieliznę średniego wieku.}
\poltrans{Time has erased the distinction between the slightly older and the 
slightly younger, throwing us all onto the sandbank of middle age, once so distant.}

\textbf{zatarł} --- perfective past of \textit{zatrzeć} (to erase, to obliterate) 
--- a total erasure, not a diminishing.

\textbf{wyrzucając nas wszystkich} --- \textit{throwing us all} --- present 
adverbial participle: not guiding but \textit{throwing}, as a wave throws 
something onto a beach.

\textbf{mielizna} --- \textit{sandbank, shoal} --- a nautical word for shallow 
water where boats run aground. Not a destination chosen but a place where you 
strand. \textbf{Tak kiedyś odległą}: the thing that seemed impossibly far 
away turned out to be right here.

% ────────────────────────────────────────────────────────────

\poltext{Dalej dotarł tylko Fiszer.}
\poltrans{Only Fiszer made it further.}

\textbf{Dalej} --- \textit{further, beyond} --- past the sandbank, which is 
to say, into death. Subject last, the heaviest word at the end. The sandbank 
metaphor inverted in four words: those stranded on the shallows are the 
fortunate ones. Perhaps the finest sentence in the chapter.

% ────────────────────────────────────────────────────────────

\poltext{Gdzie są te niegdysiejsze loki, smukłości i gładkości?!}
\poltrans{Where are those erstwhile curls, slender figures, and smooth complexions?!}

\textbf{niegdysiejsze} --- \textit{erstwhile, of former times} --- archaic and 
literary, from \textit{niegdyś} (once, in days gone by). Applied to curls and 
smooth skin at a funeral: simultaneously grandiose and comic.

The ?! combines question and cry. After the quiet devastation of 
\textit{Dalej dotarł tylko Fiszer}, the author suddenly throws his hands up.

% ────────────────────────────────────────────────────────────

\poltext{Dziewczynki i chłopcy zmienili się najpierw w panie i panów, by potem 
przeistoczyć się w\ldots kogoś takiego.}
\poltrans{The little girls and boys turned first into ladies and gentlemen, 
only to then transform into\ldots someone like this.}

\textbf{przeistoczyć się} --- stronger than \textit{zmienić się}: the prefix 
\textbf{prze-} signals a thorough, complete transformation, passing through 
into something else entirely.

\textbf{w\ldots kogoś takiego} --- the ellipsis and vague demonstrative do 
everything. The first transformation is named with dignity; the second trails 
into a gesture at the actual bodies at the actual funeral. The author cannot 
find a word for what they have become.

% ────────────────────────────────────────────────────────────

\poltext{Tylko po co tyle fałdek, zakoli i podbródków?}
\poltrans{But why so many folds, bald patches, and double chins?}

\textbf{Tylko po co} --- \textit{but why, what's the point of} --- colloquial 
and resigned, more weary than \textit{dlaczego}.

\textbf{zakoli} --- genitive plural of \textit{zakole}: primarily a river bend, 
used colloquially for receding hairlines. The geographical metaphor is quietly apt.

The diminutives (\textit{fałdek}, \textit{zakoli}, \textit{podbródków}) give 
the inventory rueful tenderness rather than pure horror.

% ────────────────────────────────────────────────────────────

\poltext{Protez, licówek, okularów?}
\poltrans{Of dentures, veneers, glasses?}

The list continues and accelerates: from soft diminutives into harder material. 
Dental hardware, cosmetic repairs, corrective lenses. The body is not just 
ageing but being actively maintained, patched, propped up.

% ────────────────────────────────────────────────────────────

\poltext{Jeżeli oni pomyśleli o mnie to, co ja o nich, to są z nich świnie!}
\poltrans{If they thought about me what I thought about them, then they are swine!}

\textbf{to są z nich świnie} --- \textit{then they are swine} --- colloquial 
construction, literally \textit{there are swine from them}. The logic is 
airtight and entirely self-implicating: the \textit{świnie} in question is 
inescapably also the author.

% ────────────────────────────────────────────────────────────

\poltext{Po chwili oczywiście nam przeszło, ale pierwsze wrażenie\ldots}
\poltrans{After a moment it of course passed, but the first impression\ldots}

\textbf{chwila} --- one of the most common time words in Polish: 
\textit{za chwilę} (in a moment), \textit{przez chwilę} (for a moment), 
\textit{po chwili} (after a moment).

\textbf{nam przeszło} --- impersonal construction, the dative \textbf{nam} 
indicating who experienced the passing. Not \textit{we stopped feeling it} 
but \textit{it passed for us}.

The ellipsis refuses to complete the thought because the thought completes 
itself. The three dots are funnier than any conclusion could be.

% ── Recognition ─────────────────────────────────────────────
\section{Rozpoznanie / Recognition}

\poltext{Hanka? Waldi? To naprawdę wy? A ta starsza pani to kto? Nie, pomyłka, 
ta pani to nie od nas --- z innego pogrzebu.}
\poltrans{Hanka? Waldi? Is it really you? And who is that older lady? No, my 
mistake, that lady isn't one of ours --- she's from a different funeral.}

\textbf{ta pani to nie od nas} --- \textit{that lady isn't one of ours} --- 
\textbf{od nas} literally \textit{from us}, meaning \textit{one of our group}.

\textbf{z innego pogrzebu} --- \textit{from a different funeral} --- the 
chapter title (\textit{Jeszcze jeden pogrzeb}) lands quietly here. They 
have aged so much that a classmate is momentarily indistinguishable from 
the elderly mourners at the next funeral over.

% ────────────────────────────────────────────────────────────

\poltext{--- No, ciebie to już dawno nie widziałam. Bardzo dawno\ldots --- 
Hanka przyglądała mi się badawczo.}
\poltrans{`Well, I haven't seen you in a long time. A very long time\ldots' 
--- Hanka was scrutinising me searchingly.}

\textbf{No} --- untranslatable particle: somewhere between \textit{well}, 
\textit{so}, and \textit{right then}. Ubiquitous in Polish conversation.

\textbf{widziałam} --- feminine ending confirms it is Hanka speaking.

\textbf{przyglądała mi się badawczo} --- \textit{was scrutinising me searchingly}. 
\textbf{Przyglądać się}: imperfective, a sustained deliberate examination. 
\textbf{Badawczo} --- from \textbf{badać} (to examine, to investigate) --- 
slightly clinical, investigative. The scrutiny is mutual: the author has been 
inventorying everyone's deterioration and is now on the receiving end of exactly 
the same examination.

% ────────────────────────────────────────────────────────────

\poltext{--- Cześć. Chciałaś może powiedzieć, że nic się nie zmieniłem?}
\poltrans{`Hi. Did you perhaps mean to say that I haven't changed at all?'}

\textbf{Chciałaś może powiedzieć} --- \textbf{chciałaś}: past feminine of 
\textit{chcieć}, confirming he addresses Hanka. \textbf{Może} softens the 
question into a tease rather than a demand.

\textbf{nic się nie zmieniłem} --- double negation (\textit{nic\ldots nie}) 
as Polish requires. He is fishing for a compliment, and he knows it. After 
Hanka's long \textit{badawczo} scrutiny he pre-empts whatever she was about 
to say by supplying the most flattering possible interpretation. The gap 
between what he hopes she might say and what her gaze actually found is 
the whole comedy of the moment.

% ────────────────────────────────────────────────────────────

\poltext{--- Nie. Nie chciałam.}
\poltrans{`No. I didn't mean to say that.'}

Two words, then two words. Hanka doesn't elaborate, doesn't soften it,
doesn't laugh. She just answers the question he asked.

\textbf{Nie chciałam} --- feminine past of \textit{chcieć} (to want).
The \textit{-am} ending confirms unambiguously that Hanka is speaking,
and makes the denial entirely personal and direct. Not \textit{nie,
pomyłka} (no, my mistake) or \textit{no, trochę się zmieniłeś} (well,
you've changed a little) --- just the flat, complete negation of the
thing he hoped she'd say.

The structure mirrors his line exactly: he asked \textit{chciałaś może
powiedzieć} --- she answers \textit{nie chciałam}. He handed her the
script; she handed it back. The comedy is in the ruthlessness of the
economy.

% ────────────────────────────────────────────────────────────

\poltext{--- To teraz będziesz widywać mnie częściej. Wybieram się do 
ciebie po szafy i kuchnię. Urządzam się --- oznajmiłem Hance, która była 
lokalną bizneswoman w branży meblowej.}
\poltrans{`So now you'll be seeing me more often. I'm heading to you for 
wardrobes and a kitchen. I'm setting myself up,' I informed Hanka, who 
was a local businesswoman in the furniture trade.}

\textbf{To} --- the conversational pivot particle, not the pronoun
\textit{to} (that). Equivalent to \textit{so then} or \textit{well
then} --- he takes Hanka's brutal honesty and immediately turns it into
an opportunity.

\textbf{będziesz widywać} --- future tense with imperfective infinitive
\textit{widywać}. Not \textit{widzieć} (to see, a single instance) but
\textit{widywać} --- the habitual/frequentative imperfective, implying
regular, recurring sightings. He is not saying \textit{you'll see me};
he is saying \textit{you'll be seeing me, regularly, as a pattern.}

\textbf{częściej} --- comparative of \textit{często} (often):
\textit{more often}. Confirms the habitual sense of \textit{widywać}.

\textbf{Wybieram się} --- \textit{I'm heading, I'm making my way}. From
\textit{wybierać się} --- to be setting out, to be planning a trip.
Slightly effortful or deliberate; not just \textit{idę} (I'm going) but
\textit{I'm making the trip to you}.

\textbf{po szafy i kuchnię} --- \textit{for wardrobes and a kitchen}.
This is the \textit{po} of purpose: the preposition meaning \textit{in
order to fetch/get}. You go \textit{po chleb} to the bakery,
\textit{po dzieci} to the school. Here: to collect wardrobes and a
kitchen. The items are named with breezy confidence, as though Hanka's
shop is simply his personal furniture depot.

\textbf{Urządzam się} --- \textit{I'm setting up / furnishing / sorting
myself out}. From \textit{urządzić się} --- to get oneself installed,
to set up house. The reflexive \textit{się} makes it personal and
ongoing.

\textbf{oznajmiłem} --- \textit{I informed, I announced}. From
\textit{oznajmić} --- a formal, declarative verb. Not
\textit{powiedziałem} (I said) but something closer to \textit{I
notified, I made known}. Applied to furniture shopping plans announced
at a funeral, it has the same comic register as \textit{z takim
postanowieniem udaliśmy się na pogrzeb} earlier: language of civic
proclamation wrapped around mundane errands.

\textbf{w branży meblowej} --- \textit{in the furniture trade/sector}.
\textit{Branża} (branch, sector, industry) is a thoroughly modern
business word. The narrator explains in passing exactly why he pivoted
from \textit{Nie chciałam} to furniture shopping in under one second.

% ── Kolejny kolega / Another schoolmate ─────────────────────
\section{Kolejny kolega / Another Schoolmate}

\poltext{Witamy się z kolejnym kolegą ze szkoły. To szef Daisy w 
przychodni stomatologicznej. Chodzili razem do klasy.}
\poltrans{We greet another schoolmate. This is Daisy's boss at a dental 
clinic. They were in the same class together.}

\textbf{Witamy się} --- \textit{we greet each other} --- first person
plural reflexive, the same \textit{my} (we) as throughout: the author
and his old classmates as a collective.

\textbf{kolejnym} --- \textit{another, the next one} --- from
\textit{kolej} (turn, sequence): yet one more in an arriving series.
The funeral keeps filling up.

\textbf{przychodnia stomatologiczna} --- \textit{dental clinic}.
\textit{Przychodnia} is an outpatient clinic or medical practice.
\textit{Stomatologiczna} --- adjectival form of \textit{stomatologia}
(dentistry), the formal Polish word, from Greek roots.

\textbf{Chodzili razem do klasy} --- \textit{they used to go to class
together}. \textbf{Chodzić do klasy} --- literally \textit{to walk to
class} --- is the standard expression for attending school together.
The imperfective past \textit{chodzili} gives it its habitual,
long-ago quality: not a single event but an entire shared childhood,
compressed into one verb.

% ────────────────────────────────────────────────────────────

\poltext{--- Cześć, Waldi. Co ci się stało w nogę?}
\poltrans{`Hi, Waldi. What happened to your leg?'}

\textbf{Co ci się stało} --- \textit{what happened to you}. The dative
\textbf{ci} (\textit{to you}) makes it personal: not \textit{what
happened} in the abstract but \textit{what befell you}. Polish uses
this dative-of-the-affected-person naturally where English simply says
``what happened to your leg'' --- but the Polish encodes a shade of
sympathy or concern that English elides.

\textbf{w nogę} --- accusative of \textit{noga} after \textit{w} ---
the accusative with \textit{w} expresses direction into something,
motion toward a target. Not \textit{w nodze} (locative: \textit{in the
leg}, where something is located) but \textit{w nogę} --- something
went \textit{into} the leg, struck it. The construction is idiomatic
for injuries: \textit{dostał w nogę} --- \textit{he got it in the leg}.

The question tells us everything. After all the inventory of folds and
receding hairlines and dental hardware, here is something more acute:
one of them is visibly injured. The question is asked directly, without
preamble, the way old friends ask things.

% ────────────────────────────────────────────────────────────

\poltext{--- Chłop z kryką? Nie za wcześnie? --- dodaje Groszek.}
\poltrans{`A bloke with a crutch? Isn't it a bit early for that?' --- 
Groszek adds.}

\textbf{Chłop} --- \textit{bloke, fellow, man}. Colloquial and earthy
--- not the neutral \textit{mężczyzna} but something closer to
\textit{guy, chap}. Here it has a gruff, joshing quality: \textit{a
grown man, on a crutch?}

\textbf{kryką} --- instrumental of \textit{kryka} --- \textit{crutch},
specifically an underarm crutch. The bluntness of the phrasing ---
subject-plus-instrument, no verb --- reads like a double-take.

\textbf{Nie za wcześnie?} --- \textit{Isn't it a bit early? Too soon?}
\textbf{Za} is the \textit{too} of excess: \textit{za dużo} (too
much), \textit{za późno} (too late), \textit{za wcześnie} (too early).
The affectionate cruelty is precise: you're not old enough yet to need
that. The implication being, of course, that the day will come for all
of them.

\textbf{Groszek} --- literally \textit{little pea}: a childhood
nickname, carried forward unchanged. The diminutive suffix
\textit{-ek} is everywhere in Polish nicknames; \textit{groch} (pea)
becomes \textit{groszek}. That this person is still called Little Pea
at a funeral in middle age is quietly perfect.

% ── Serwis pogwarancyjny / Post-warranty Servicing ──────────
\section{Serwis pogwarancyjny / Post-warranty Servicing}

\poltext{--- Endoproteza. W tym wieku jest już potrzebny serwis 
pogwarancyjny. Niedawno była operacja, jeszcze trochę pochodzę o kulach.}
\poltrans{`A joint replacement. At this age you already need post-warranty 
servicing. There was an operation recently, I'll be walking on crutches a 
little while longer.'}

\textbf{Endoproteza} --- \textit{endoprosthesis, joint replacement}: an
artificial joint, typically hip or knee. Dropped alone, without a verb,
as a complete explanation. One word, clinical, total.

\textbf{W tym wieku} --- \textit{at this age} --- echoing Hanka's
\textit{w twoim wieku} from moments earlier. The phrase is becoming a
quiet refrain throughout this gathering.

\textbf{serwis pogwarancyjny} --- \textit{post-warranty servicing}: the
maintenance required once the manufacturer's guarantee has expired.
Entirely a term from consumer electronics and appliances. Waldi has
reframed his own body as a piece of machinery that has outlived its
warranty period and now requires paid maintenance. \textbf{Już} adds
the familiar rueful note: the milestone has been reached, perhaps
sooner than expected.

\textbf{Niedawno była operacja} --- \textit{there was an operation
recently}. Impersonal construction: not \textit{I had an operation} but
\textit{an operation was} --- the event stated as a fact, Waldi at a
slight grammatical remove from his own surgery.

\textbf{jeszcze trochę pochodzę o kulach} --- \textit{I'll walk on
crutches a little while longer}. \textbf{Pochodzę} --- future with the
prefix \textit{po-} giving a sense of duration covering a period:
\textit{I'll be going about} for a while. \textbf{O kulach} ---
\textit{on crutches}: instrumental of \textit{kule}. \textbf{Jeszcze
trochę} --- \textit{just a bit more, a little while yet} --- patient
and matter-of-fact.

The \textit{serwis pogwarancyjny} metaphor is the centrepiece. It
extends naturally from the whole chapter's inventory of bodily
deterioration --- folds, veneers, glasses, dentures --- but Waldi
crystallises it into a single perfectly chosen phrase. What makes it
land is that he says it without complaint, almost cheerfully: he has
accepted the terms of the post-warranty period and is simply reporting
on current service status.

% ────────────────────────────────────────────────────────────

\poltext{--- Do wesela się zagoi --- dodaje optymistycznie Daisy.}
\poltrans{`It'll heal by the wedding' --- Daisy adds optimistically.}

\textbf{Do wesela się zagoi} --- a set Polish folk reassurance:
\textit{it'll heal by the wedding}, meaning \textit{don't worry, you'll
be fine in time for what matters}. Mothers say it to children who have
scraped their knees. Applied to a man in his sixties with a hip
replacement, standing at a funeral, it is either the most cheerful
possible response --- or, depending on your angle, quietly devastating.
At this age, \textit{do wesela} might mean a grandchild's wedding,
still years away.

\textbf{zagoi się} --- from \textit{zagoić się}, to heal over. The
\textit{za-} prefix signals completion: fully healed, closed over.
Future perfective.

\textbf{dodaje optymistycznie Daisy} --- the adverb
\textbf{optymistycznie} is placed by the narrator with a certain
dryness: he notes the optimism as a quality of the remark, not
necessarily endorsing it. Three different responses to the same crutch:
Groszek mocks (\textit{nie za wcześnie?}), Waldi explains with rueful
wit (\textit{serwis pogwarancyjny}), Daisy closes with breezy folk
reassurance. Three temperaments, a handful of lines, a graveside.

% ────────────────────────────────────────────────────────────

\poltext{--- No, najpierw to musiałby się rozwieść. Polecam, ja tam 
bardzo dobrze na tym wyszłam --- w ten sposób podsumowuje Hanka ostatnie 
rewelacje w swoim życiu.}
\poltrans{`Well, first he'd have to get divorced. I recommend it --- I 
came out of it very well' --- that's how Hanka summarises the latest 
revelations in her life.}

\textbf{najpierw to musiałby się rozwieść} --- \textit{first he'd have
to get divorced}. Hanka has understood the furniture shopping not as a
business opportunity but as the announcement of a new living situation,
and has drawn the correct inference. \textbf{Musiałby} --- conditional
of \textit{musieć}: \textit{he would have to} --- the condition stated
calmly, as a logical preliminary. \textbf{Rozwieść się} ---
\textit{to get divorced}, perfective reflexive.

\textbf{Polecam} --- \textit{I recommend it.} First person singular,
present, no hedging. Not \textit{it wasn't so bad} or \textit{it has
its advantages} --- a clean, direct endorsement. The same word you
would use recommending a restaurant.

\textbf{ja tam bardzo dobrze na tym wyszłam} --- \textit{I came out of
it very well}. \textbf{Tam} here is not spatial but a colloquial
particle adding casual personal emphasis: \textit{for my part, as for
me}. \textbf{Wyjść na czymś} --- \textit{to come out of something} in
terms of result: to fare well or badly. \textbf{Bardzo dobrze} ---
Hanka is not merely surviving her divorce; she came out ahead. The
furniture business, presumably, is part of the evidence.

\textbf{ostatnie rewelacje w swoim życiu} --- \textit{the latest
revelations in her life}. \textbf{Rewelacje} carries a flavour of
\textit{bombshells, sensational news} --- the word you'd use for gossip
or surprising disclosures. Applied to one's own divorce, stated
cheerfully at a funeral, it has a wonderfully self-aware theatricality.
The narrator labels it all a \textit{summary}, watching Hanka perform
her own life story with total composure.

In two sentences Hanka has: correctly diagnosed the author's domestic
situation, volunteered her own experience as precedent, rendered a
verdict (\textit{polecam}), reported her outcome (\textit{bardzo
dobrze}), and had it all dryly labelled by the narrator. Not a word
wasted.

% ────────────────────────────────────────────────────────────

\poltext{--- Chodźcie do Fiszera. Potem sobie pogadacie --- postanawiam.}
\poltrans{`Go to Fiszer. You can chat afterwards' --- I decide.}

\textbf{Chodźcie do Fiszera} --- \textit{go to Fiszer / go to Fiszer's
grave}. \textbf{Chodźcie} is the second person plural imperative of
\textit{chodzić}: \textit{go, come on, move}. The group has been
standing around exchanging divorces and hip replacement diagnoses; now
the author redirects them. \textit{Do Fiszera}: genitive, the direction
toward Fiszer. Not \textit{to the grave}, not \textit{to the ceremony}
--- just \textit{to Fiszer}, as if he's waiting nearby and they've kept
him long enough.

\textbf{Potem sobie pogadacie} --- \textit{you can chat afterwards}.
\textbf{Sobie} is the reflexive dative used as a casual particle ---
it softens and domesticates the verb, implying something done at
leisure, for one's own pleasure: \textit{have yourselves a chat}.
\textbf{Pogadać} is the perfective of the colloquial \textit{gadać}
(to chat, to natter) --- not the formal \textit{rozmawiać}.

\textbf{postanawiam} --- \textit{I decide, I resolve}. Present tense,
first person singular. Not \textit{I said} or \textit{I suggested} but
\textit{I decide} --- the same weighty \textit{postanowienie}
(resolution) from earlier in the chapter, now in verb form.

All the reunion warmth has been real, but \textit{chodźcie do Fiszera}
brings everything back. They are, after all, at a funeral. First things
first.

% ── Nieobecni / The Absent ───────────────────────────────────
\section{Nieobecni / The Absent}

\poltext{--- A nieobecni? Sławek nie przyjdzie. Jest znowu gdzieś na 
końcu świata --- wyjaśnia siostra.}
\poltrans{`And those who aren't here? Sławek won't come. He's somewhere 
at the end of the world again' --- his sister explains.}

\textbf{A nieobecni?} --- \textit{And those who aren't here?}
\textbf{Nieobecni} --- \textit{the absent ones}, plural adjective used
as a noun, from \textit{obecny} (present) with the negating
\textit{nie-}. Someone takes a mental roll call. The funeral is also
an accounting of who didn't make it.

\textbf{Jest znowu gdzieś na końcu świata} --- \textit{he's somewhere
at the end of the world again}. \textbf{Znowu} --- \textit{again} ---
the word that carries the whole character sketch. This is not new
behaviour; Sławomir is habitually, reliably somewhere unreachable.
\textbf{Gdzieś} --- \textit{somewhere} --- the vagueness deliberate.
\textbf{Na końcu świata} --- \textit{at the end of the world} --- a
fixed Polish expression for an extremely remote place, equivalent to
\textit{the back of beyond}.

\textbf{wyjaśnia siostra} --- \textit{his sister explains}. She has
clearly explained Sławek's whereabouts before --- perhaps many times.
The verb \textit{wyjaśniać} implies there was something requiring
clarification; she supplies the answer to the question the group was
already asking.

% ────────────────────────────────────────────────────────────

\poltext{Sławomir był artystą konceptualnym (jak mawia Daisy, 
współczesna sztuka nie jest od tego, by ją rozumieć), wykładowcą 
i podróżnikiem, postacią nieprzeciętną.}
\poltrans{Sławomir was a conceptual artist (as Daisy likes to say, 
contemporary art isn't there to be understood), a lecturer and 
traveller, an exceptional figure.}

\textbf{Sławomir} --- the full name, now that we're in the narrator's
expository voice rather than dialogue. He was \textit{Sławek} to the
group; here he gets the dignity of his proper name.

\textbf{był artystą konceptualnym} --- \textit{was a conceptual
artist}. Instrumental after \textit{być} (to be) in the past:
\textbf{artystą}, instrumental of \textit{artysta}. Polish requires
this throughout --- you don't \textit{be} a thing in the nominative
but \textit{be} it in the instrumental, as though identity is something
you inhabit rather than something you simply are.

\textbf{jak mawia Daisy} --- \textit{as Daisy likes to say}.
\textbf{Mawia} is a frequentative form of \textit{mówić} --- not
\textit{says} on this occasion but \textit{is wont to say, habitually
says}. It marks the parenthesis as a well-worn Daisy opinion, trotted
out reliably.

\textbf{współczesna sztuka nie jest od tego, by ją rozumieć} ---
\textit{contemporary art isn't there to be understood}. \textbf{Od
tego, by} --- \textit{for that purpose, in order to} --- a
construction meaning something exists (or doesn't exist) for a given
purpose. Daisy's parenthetical is both a defence of Sławomir's work
and a perfectly self-contained piece of comic deflation.

\textbf{podróżnik} --- a wonderfully old-fashioned, romantic word for
traveller, explorer: not a \textit{turysta} (tourist) but someone for
whom travel is a vocation or identity.

\textbf{postacią nieprzeciętną} --- \textit{an exceptional figure}.
\textbf{Nieprzeciętny} --- literally \textit{not average}, from
\textit{przeciętny} (average, mediocre). The summary epithet, placed
last with the weight of a conclusion.

The parenthetical is exquisite: rather than explaining conceptual art,
the narrator outsources the explanation to Daisy --- who doesn't
explain it either, but dismisses the need for explanation altogether.
Sławomir's work is described, defended, and gently mocked in eleven
words, none of which actually tell us anything about it.

% ────────────────────────────────────────────────────────────

\poltext{Mało małomiasteczkową.}
\poltrans{Not very small-town.}

\textbf{Mało} --- \textit{not very, scarcely} --- used as a negating
qualifier: \textit{hardly, not much of a}. It doesn't flatly negate
but diminishes.

\textbf{małomiasteczkową} --- instrumental feminine of
\textbf{małomiasteczkowy}: \textit{small-town, provincial}. A compound
adjective built from \textbf{małe miasto} (small town). It agrees
instrumentally with \textit{postacią} from the previous sentence.

The sentence functions as a dry afterthought: having assembled
Sławomir's credentials --- artist, lecturer, traveller, exceptional ---
the narrator adds this fragment as if \textit{małomiasteczkowy} were
the natural opposite of everything just listed, and confirming he
wasn't that is the final, clinching point. There is also a faint
self-awareness: these are people \textit{from} a small town, gathered
at a funeral in it. Sławomir escaped.

% ────────────────────────────────────────────────────────────

\poltext{Miał być, ale ostatecznie nie dojechał.}
\poltrans{He was supposed to come, but in the end he didn't make it.}

\textbf{Miał być} --- \textit{he was supposed to be there, he was
meant to come}. \textbf{Mieć} + infinitive is the Polish construction
for obligation, expectation, or plan: \textit{miał przyjść} --- he was
supposed to come; \textit{miał być} --- he was supposed to be there.
It implies an arrangement that existed, a slot reserved for him.

\textbf{ostatecznie nie dojechał} --- \textit{in the end he didn't
make it}. \textbf{Ostatecznie} --- \textit{ultimately, in the end} ---
carries a slight air of resignation, of a situation that played out to
its conclusion. \textbf{Dojechać} --- perfective of \textit{dojeżdżać},
to arrive by vehicle, to complete a journey. The \textit{do-} prefix
is directional completion: to travel \textit{all the way to} a
destination. Not \textit{nie przyszedł} (didn't come) but \textit{nie
dojechał} --- the gap between setting out and arriving, the journey
that didn't finish.

% ────────────────────────────────────────────────────────────

\poltext{Zabrakło też Mirka, filozofa amatora i diamentowego 
konsultanta Amwaya.}
\poltrans{Mirek was also absent --- amateur philosopher and diamond 
consultant for Amway.}

\textbf{Zabrakło} --- impersonal perfective of \textit{brakować} (to
lack, to be missing). Not \textit{Mirek nie przyszedł} (Mirek didn't
come) but \textit{there was a lacking of Mirek} --- the person rendered
as an absence, a gap in the room. \textbf{Mirka}: genitive, as required
by the impersonal \textit{zabrakło}.

\textbf{filozofa amatora} --- \textit{amateur philosopher}. Both in
the genitive, in apposition to \textit{Mirka}. \textbf{Amator} in
Polish carries less condescension than in English; it means someone
who does something from love (\textit{amor}) rather than
professionally.

\textbf{diamentowego konsultanta Amwaya} --- \textit{diamond consultant
for Amway}. \textbf{Diamentowy} --- \textit{diamond} --- is an actual
rank in Amway's multi-level marketing hierarchy, awarded to top-tier
distributors. The collision of \textit{filozofa amatora} and
\textit{diamentowego konsultanta Amwaya} in a single breath is the
entire joke: Mirek has two identities, and the narrator presents them
with perfect deadpan equivalence. Neither is privileged over the other.

% ────────────────────────────────────────────────────────────

\poltext{Był nieobecnym usprawiedliwionym, na stałe mieszkał we 
Włoszech. Łączył się w bólu na Facebooku.}
\poltrans{He was an excused absence --- he lived permanently in Italy. 
He joined in the grief on Facebook.}

\textbf{nieobecnym usprawiedliwionym} --- \textit{an excused absence}.
In Polish schools, absences are formally categorised as
\textbf{usprawiedliwione} (excused, with a valid reason) or
\textbf{nieusprawiedliwione} (unexcused). The narrator applies the
bureaucratic vocabulary of a primary school attendance book to a grown
man's failure to attend his old classmate's funeral. Since Mirek lives
in Italy, his absence is officially \textit{usprawiedliwione}.

\textbf{na stałe} --- \textit{permanently, for good}. Not
\textit{currently} or \textit{for now} but settled, fixed.

\textbf{Łączył się w bólu na Facebooku} --- \textit{he joined in the
grief on Facebook}. \textbf{Łączyć się w bólu} --- \textit{to unite
in pain/grief} --- is the formal Polish condolence phrase used in
funeral notices and obituaries: \textit{I share in your grief, I unite
with you in sorrow}. The ancient, solemn phrase has been redirected
through a social media platform. The narrator records it without
editorial comment. None is needed.

% ── Umarła klasa / The Dead Class ───────────────────────────
\section{Umarła klasa / The Dead Class}

\poltext{A martwi? Po uroczystości Groszek, czyniąc cmentarne honory 
domu, wyciągnął wreszcie okolicznościową flaszkę i zaproponował, by 
odwiedzić resztę naszych na kwaterach, tę „umarłą klasę".}
\poltrans{And the dead? After the ceremony, Groszek, playing the 
graveyard host, finally produced the occasion's bottle and proposed 
that we visit the rest of our people in their plots --- that `dead class'.}

\textbf{A martwi?} --- \textit{And the dead?} Two words, and the
question turns the whole register on its head. We've had the absent
living --- Sławomir at the end of the world, Mirek on Facebook from
Italy. Now the narrator pivots to the other category of non-attenders:
those in the ground.

\textbf{czyniąc cmentarne honory domu} --- \textit{playing the
graveyard host}, literally \textit{doing the cemetery honours of the
house}. \textbf{Czyniąc} --- present adverbial participle of
\textit{czynić} (to do, to make) --- slightly elevated, literary.
\textbf{Honory domu} --- \textit{the honours of the house} --- a fixed
phrase for the duties of a host, transplanted wholesale to a cemetery,
so that Groszek becomes the host of this outdoor gathering among the
graves.

\textbf{wyciągnął wreszcie okolicznościową flaszkę} --- \textit{finally
produced the occasion's bottle}. \textbf{Wyciągnął} --- \textit{pulled
out, produced} --- the perfective suggesting the bottle has been
waiting, held in reserve. \textbf{Wreszcie} --- \textit{finally, at
last}. \textbf{Okolicznościową} --- \textit{for the occasion,
commemorative}: an adjective applied to things made or brought for a
specific event. This is the fulfilment of \textit{zdecydowaliśmy, że
dla uczczenia jego pamięci trzeba się napić} from the very opening of
the chapter.

\textbf{na kwaterach} --- \textit{in their plots}. \textbf{Kwatera}
means a plot or section in a cemetery, but also a billet or lodging
(military quarters, accommodation). Visiting the dead in their
\textit{kwatery} as one might visit friends in their lodgings.

\textbf{tę „umarłą klasę"} --- \textit{that `dead class'}. A direct
reference to Tadeusz Kantor's legendary 1975 theatre piece
\textit{Umarła klasa}, one of the defining works of Polish
avant-garde theatre, in which elderly figures return to their school
desks and encounter the ghosts of their childhood selves. The narrator
names it quietly, without explanation, trusting the reader. The
classmates gathering at Fiszer's grave, the absent ones already under
it --- they are Kantor's image made literal.

% ────────────────────────────────────────────────────────────

\poltext{Oczywiście w tym momencie zmyślam. Groszek by tak nie 
powiedział. Groszek był po politechnice.}
\poltrans{Obviously at this point I'm making things up. Groszek 
wouldn't have said it like that. Groszek had been to polytechnic.}

\textbf{zmyślam} --- present tense of \textit{zmyślać} (to make up,
to fabricate, to invent) --- not apologetic, not alarmed, just stating
a fact. The \textit{oczywiście} (obviously) is the masterstroke: of
course he's inventing. Did you think he wasn't?

\textbf{Groszek by tak nie powiedział} --- \textit{Groszek wouldn't
have said it like that}. The narrator has caught himself putting
\textit{cmentarne honory domu} and \textit{tę umarłą klasę} into
Groszek's mouth --- language that is clearly the author's, literary
and allusive --- and corrects the record.

\textbf{Groszek był po politechnice} --- \textit{Groszek had been to
polytechnic}. The construction \textit{być po [instytucji]} means to
have attended and graduated: \textit{po wojsku} (having done military
service), \textit{po studiach} (having finished university). The
polytechnic --- a technical university, engineering-focused --- is
offered as the complete explanation. Engineers do not say \textit{tę
umarłą klasę}.

The joke runs in two directions: gently at Groszek's expense (engineers
don't talk like that) and at the author's own (I was showing off and I
know it). The self-correction is funnier than the original flourish,
and the original flourish was already very good.

% ────────────────────────────────────────────────────────────

\poltext{Są tu prawie wszyscy nasi klasowi umarli! Nie chce się 
wierzyć! Ze łzami w oczach pociągnęliśmy z gwinta i ruszyliśmy 
w obchód.}
\poltrans{Almost all our class dead are here! It's hard to believe! 
With tears in our eyes we took a pull from the bottle and set off 
on our rounds.}

\textbf{Są tu prawie wszyscy nasi klasowi umarli!} --- what Groszek
actually said. Plain, direct, with an exclamation mark.
\textbf{Klasowi umarli} --- \textit{class dead}, the dead of the class
--- a noun phrase of remarkable bluntness. Not \textit{ci, których
straciliśmy} (those we have lost) but simply \textit{umarli} --- the
dead ones, \textit{ours}. \textbf{Prawie wszyscy} --- \textit{almost
all} --- the \textit{prawie} doing quiet, terrible work: not all yet,
but nearly.

\textbf{Nie chce się wierzyć!} --- \textit{It's hard to believe! One
doesn't want to believe it!} The impersonal \textit{chce się}
construction: not \textit{I don't want to believe it} but \textit{it
doesn't want to be believed}, \textit{belief resists}.

\textbf{pociągnęliśmy z gwinta} --- \textit{we took a pull from the
bottle}. \textbf{Pociągnąć z gwinta} is a vivid colloquial expression:
\textbf{gwint} literally means a screw-thread --- the neck of the
bottle where the cap screws on. To drink \textit{z gwinta} is to drink
straight from the bottle, no glasses.

\textbf{ruszyliśmy w obchód} --- \textit{we set off on our rounds}.
\textbf{Obchód} --- a round, a circuit, an inspection tour: the word
used for a guard's patrol, a doctor's ward round, a ceremonial
procession. Here: a tour of the graves of their dead classmates, bottle
in hand, tears in eyes, Groszek leading the way.

The chapter has been moving toward this moment from its first sentence.
The resolution to drink is honoured; Groszek's plain words --- not the
narrator's literary flourishes --- name what they are doing and why.

% ── Obchód / The Rounds ──────────────────────────────────────
\section{Obchód / The Rounds}

\poltext{Zostawiliśmy świeży grób Fiszera z wieńcami, kwiatami.}
\poltrans{We left Fiszer's fresh grave with its wreaths, its flowers.}

\textbf{Zostawiliśmy} --- \textit{we left, we left behind}. Perfective
past of \textit{zostawiać} --- to leave something in place, to leave
it as it is. Not \textit{opuściliśmy} (we departed from) but
\textit{zostawiliśmy} --- they left it there, behind them, as they
moved on.

\textbf{świeży grób} --- \textit{the fresh grave}. \textbf{Świeży}
--- \textit{fresh} --- the same word for fresh bread, fresh air,
fresh paint. Applied to a grave it is entirely ordinary Polish usage,
but the word still lands with its full weight: the earth newly turned,
not yet settled.

\textbf{z wieńcami, kwiatami} --- \textit{with its wreaths, its
flowers}. The comma between them, no \textit{i} (and), gives the list
an almost visual quality --- the eye moving across the grave's surface,
noting one thing, then another.

They set off on their rounds --- and the first thing they do is leave
Fiszer behind. The living move on; the fresh grave stays.

% ────────────────────────────────────────────────────────────

\poltext{Z przyjmującą kondolencje zapłakaną i śmiertelnie zmęczoną żoną.}
\poltrans{With the tear-stained and mortally exhausted wife receiving 
condolences.}

The sentence continues the list begun with \textit{z wieńcami,
kwiatami} --- the grave left behind, and now: the widow.

\textbf{Z przyjmującą kondolencje} --- \textit{with the one receiving
condolences}. \textbf{Przyjmującą} is the present active participle of
\textit{przyjmować} (to receive, to accept) in the instrumental
feminine. She is not described as standing, or waiting, or grieving
--- but as \textit{receiving condolences}: defined entirely by the
ritual function she is currently performing.

\textbf{zapłakaną} --- \textit{tear-stained, having been crying}.
Instrumental feminine participle from \textit{zapłakać się} --- to
have cried, to be red-eyed with weeping. The \textit{za-} prefix marks
the completed state: she has been crying and it shows.

\textbf{śmiertelnie zmęczoną} --- \textit{mortally exhausted}.
\textbf{Śmiertelnie} --- \textit{mortally, deathly} --- the intensifier
drawn from death itself, entirely idiomatic in Polish
(\textit{śmiertelnie znudzony} --- bored to death). Here it sits with
particular weight: she is deathly tired, at a death, having transported
her husband's body from the south of France. The logistics noted so
dryly at the chapter's opening --- \textit{musiała transportować jego
zwłoki} --- have arrived at their human cost.

% ────────────────────────────────────────────────────────────

\poltext{Niewprawnym zygzakiem poruszaliśmy się za Groszkiem pomiędzy 
grobami po wąskich cmentarnych ścieżkach, szukając klasowych 
nieboszczyków.}
\poltrans{In an unsteady zigzag we moved behind Groszek between the 
graves along the narrow cemetery paths, looking for the class dead.}

\textbf{Niewprawnym zygzakiem} --- \textit{in an unsteady zigzag}.
The instrumental of manner. \textbf{Niewprawny} --- \textit{unskilled,
unpractised, clumsy} --- from \textit{wprawny} (skilled) with the
negating \textit{nie-}. The bottle from Groszek's pocket has done its
early work. The group is weaving.

\textbf{poruszaliśmy się} --- \textit{we moved, we made our way}.
Imperfective past, the reflexive \textit{się} with \textit{poruszać}
giving the sense of sustained, ongoing movement.

\textbf{za Groszkiem} --- \textit{behind Groszek}. Instrumental after
\textit{za}. Groszek leads; the others follow. He is, after all, the
one doing the \textit{honory domu}.

\textbf{po wąskich cmentarnych ścieżkach} --- \textit{along the narrow
cemetery paths}. \textbf{Po} with the locative expresses movement
along a surface: \textit{po ulicy} (along the street), \textit{po
podłodze} (across the floor). The paths are narrow; the zigzag is
wide; the combination explains itself.

\textbf{szukając klasowych nieboszczyków} --- \textit{looking for the
class dead}. \textbf{Nieboszczyk} --- \textit{the deceased} --- a
slightly old-fashioned, almost tender word, softer than \textit{trup}
(corpse), more colloquial than \textit{zmarły} (the departed). Here
in the plural, as a category: the classmates who got there before them.

The image is both physically comic and quietly moving: a small group
weaving behind Little Pea through a Polish cemetery, bottle in hand,
searching for their dead.

% ── Cmentarny opis / Cemetery Survey ────────────────────────
\section{Cmentarny opis / Cemetery Survey}

\poltext{Tandetne walentynkowe grobowce nowych zmarłych z 
asymetrycznymi serduszkami obok zmurszałych, starych niemieckich 
pomników.}
\poltrans{Tawdry Valentine's Day tombs of the newly dead with asymmetric 
little hearts, next to crumbling, old German monuments.}

A sentence with no main verb --- a pure observation, a painter's note.
The eye sweeps the cemetery and records what it sees.

\textbf{Tandetne} --- \textit{tawdry, cheap, tacky}. From
\textit{tandeta} --- cheap goods, kitsch. An unsparing word.

\textbf{walentynkowe} --- \textit{Valentine's Day}, used as an
adjective. Applied to grave monuments: the aesthetic of Valentine's
Day --- pink, heart-shaped, decorative --- transposed onto funerary
architecture.

\textbf{z asymetrycznymi serduszkami} --- \textit{with asymmetric
little hearts}. \textbf{Serduszkami} --- diminutive instrumental plural
of \textit{serduszko}, itself a diminutive of \textit{serce} (heart).
The double diminutive: not \textit{hearts} but \textit{little hearts},
the kind on greetings cards. \textbf{Asymetrycznymi} --- the detail
that makes it real. Not even well-made little hearts.

\textbf{zmurszałych, starych niemieckich pomników} --- \textit{crumbling,
old German monuments}. \textbf{Zmurszałych} --- \textit{crumbling,
decayed, rotted} --- describing stone that has broken down with age
and damp. The cemetery predates the postwar border shifts; this part
of Silesia was German until 1945, and the old German graves remain,
weathering beside their newer Polish neighbours. The aesthetic contrast
is also a historical one.

% ────────────────────────────────────────────────────────────

\poltext{Skremowana klasa ekonomiczna w lastryko obok otoczonej 
masywnym łańcuchem, zapuszczonej mogiły Bojowników o Wolność z Dawno 
Minionych Czasów.}
\poltrans{The cremated economy class in terrazzo, next to the overgrown 
tomb of the Fighters for Freedom from Times Long Past, surrounded by 
a massive chain.}

\textbf{Skremowana klasa ekonomiczna} --- \textit{the cremated economy
class}. \textbf{Klasa ekonomiczna} --- \textit{economy class}, the
airline term for the cheapest seats, applied to those who could not
afford a full grave plot and were cremated into smaller, cheaper
monuments. The narrator has classified the dead by their funerary
budget using the language of air travel.

\textbf{w lastryko} --- \textit{in terrazzo}. \textbf{Lastryko} is
terrazzo --- the composite stone material ubiquitous in communist-era
Polish construction, common for cheaper grave surrounds. Functional,
durable, aesthetically modest.

\textbf{zapuszczonej} --- \textit{overgrown, neglected}. From
\textit{zapuścić} --- to let go, to allow to become overgrown. No one
tends this grave.

\textbf{Bojownicy o Wolność z Dawno Minionych Czasów} --- the official
communist-era designation for resistance fighters and partisans:
\textit{Fighters for Freedom} --- a bureaucratic honorific carved into
stone on countless Polish monuments. \textbf{Z Dawno Minionych Czasów}
--- \textit{from Times Long Past} --- and here the capitals are the
narrator's, not the monument's. The phrase is elevated, almost
fairy-tale in register (\textit{dawno, dawno temu} --- once upon a
time). The official heroism now left to go to seed behind a massive
chain.

% ────────────────────────────────────────────────────────────

\poltext{Sylikonowy anioł z wtryskarki obok sepiowych fotografii z 
niemodnymi krawatami i fryzurami w stylu świętej pamięci Waldemara 
Koconia.}
\poltrans{A silicone angel from an injection moulding machine, next to 
sepia photographs with unfashionable ties and hairstyles in the style 
of the late, lamented Waldemar Kocoń.}

\textbf{Sylikonowy anioł z wtryskarki} --- \textit{a silicone angel
from an injection moulding machine}. \textbf{Wtryskarki} --- genitive
of \textit{wtryskarka}, an injection moulding machine. The angel did
not come from a sculptor's hands but from a factory production line,
extruded in silicone. The word \textit{wtryskarki} contains the entire
critique.

\textbf{sepiowych fotografii} --- \textit{sepia photographs}. The oval
ceramic or glass photographs mounted on Polish graves, showing the
deceased in life --- a widespread and distinctive tradition.

\textbf{z niemodnymi krawatami i fryzurami} --- \textit{with
unfashionable ties and hairstyles}. The photographs capture their
subjects in the fashion of the decade of their death, now dated. The
ties are wide or narrow in the wrong way; the hair is set in a style
that places them precisely in time.

\textbf{w stylu świętej pamięci Waldemara Koconia} --- \textit{in the
style of the late, lamented Waldemar Kocoń}. \textbf{Świętej pamięci}
--- \textit{of blessed/holy memory} --- the traditional Polish
honorific for the deceased. Here deployed with perfect comic precision:
Waldemar Kocoń becomes the unwitting style reference for an entire era
of Polish men's grooming.

% ────────────────────────────────────────────────────────────

\poltext{Nazwiska niemieckie, polskie, nieczytelne. Cmentarny opis 
obyczajów minionego wieku.}
\poltrans{Names German, Polish, illegible. A cemetery account of the 
customs of the past century.}

\textbf{Nazwiska niemieckie, polskie, nieczytelne} --- three adjectives,
no verb. The list moves from legible history --- German, Polish, the
two peoples who have buried their dead here --- to \textit{nieczytelne}:
illegible, worn beyond reading. Time doesn't distinguish between German
and Polish; it erases them equally.

\textbf{Cmentarny opis obyczajów minionego wieku} --- \textit{a
cemetery account of the customs of the past century}. The narrator
steps back and names what this whole passage has been: a
\textit{cmentarny opis}, a survey conducted in a cemetery.
\textbf{Obyczaje} --- customs, manners, the characteristic practices
that define a time and place. The Valentine's hearts, the silicone
angel, the sepia photographs, the crumbling German stone --- all of
it the material of a social historian's survey, conducted on foot,
slightly unsteadily, with a bottle.

% ── Grób Roberta / Robert's Grave ───────────────────────────
\section{Grób Roberta / Robert's Grave}

\poltext{Po krótkich poszukiwaniach stanęliśmy obok grobu, który, 
zamiast pomnika, wieńczyła marmurowa piłka.}
\poltrans{After a brief search we stopped beside a grave which, instead 
of a headstone, was crowned by a marble ball.}

\textbf{stanęliśmy} --- \textit{we stopped, we came to a stand}.
Perfective past of \textit{stanąć}: to stop, to come to a halt, to
take up a position. The movement ceases; they have arrived.

\textbf{który, zamiast pomnika, wieńczyła marmurowa piłka} --- the
relative clause delivers its surprise with structural delay:
\textit{który} (which) --- \textit{zamiast pomnika} (instead of a
headstone) --- \textit{wieńczyła} (was crowned by) ---
\textit{marmurowa piłka} (a marble ball). The subject arrives last.
\textbf{Wieńczyć} --- \textit{to crown, to top, to cap} --- a verb
with ceremonial weight: to crown a building, to crown an achievement.
Here: to sit atop a grave. Not a cross, not an angel, not a stone slab.
A ball. In marble.

% ────────────────────────────────────────────────────────────

\poltext{--- To już druga, mniejsza. Poprzednią zajebali --- wyjaśnił 
nasz przewodnik.}
\poltrans{`This is already the second one, the smaller one. They nicked 
the previous one' --- our guide explained.}

\textbf{To już druga, mniejsza} --- \textit{this is already the second
one, the smaller one}. \textbf{Już} marks the sequence: the marble
ball is a replacement, and a downgraded one.

\textbf{Poprzednią zajebali} --- \textit{they nicked the previous one}.
\textbf{Zajebali} --- vulgar perfective past of \textit{zajebiać}, to
steal, to nick. Thoroughly common in spoken Polish, it does something
no polite synonym could: it combines theft with exasperated resignation.
\textit{Ukradli} (they stole) would be neutral; \textit{zajebali}
implies this is the kind of thing that happens to marble balls left in
cemeteries, and everyone knows it.

\textbf{nasz przewodnik} --- \textit{our guide}. \textbf{Przewodnik}
--- the same word for a tour guide, a guidebook, a scout leader.
Groszek, leading them \textit{niewprawnym zygzakiem}, is now formally
elevated to \textit{nasz przewodnik}.

% ────────────────────────────────────────────────────────────

\poltext{Grób Roberta. Adam leży obok. Adam i Robert Mitwerandu. 
Czarni śląscy bracia. Czarniejsi od węgla!}
\poltrans{Robert's grave. Adam lies beside it. Adam and Robert 
Mitwerandu. Black Silesian brothers. Blacker than coal!}

\textbf{Grób Roberta} --- two words, a new name. The \textit{obchód}
moves on.

\textbf{Adam leży obok} --- \textit{Adam lies beside it}.
\textbf{Leży} --- \textit{lies}, present tense of \textit{leżeć}:
the verb used for lying down but also for the dead in their graves.
\textit{Tu leży} --- \textit{here lies} --- appears on every headstone.
The present tense gives it quiet persistence.

\textbf{Adam i Robert Mitwerandu} --- the surname arrives and
everything shifts. Unmistakably African in origin: these are two
brothers buried side by side in a Silesian cemetery, with a name from
another continent.

\textbf{Czarni śląscy bracia} --- \textit{Black Silesian brothers}.
No verb, just the designation, stated plainly. The combination is not
a paradox to the narrator --- it is simply a fact: they were Black,
they were Silesian, they were brothers, they are here.

\textbf{Czarniejsi od węgla!} --- \textit{Blacker than coal!} The
exclamation mark carries warmth, not shock. Coal is the defining
substance of Silesia --- the region was built on it, its identity is
inseparable from it. To be \textit{czarniejsi od węgla} in Silesia
is to be measured against the local absolute. The voice is that of
someone who knew them, saying this with affection and memory.

% ────────────────────────────────────────────────────────────

\poltext{Owoce internacjonalistyczno-seksualnych uniesień czarnoskórego 
studenta, który po skończeniu w Polsce studiów zabrał się z powrotem 
na Czarny Ląd do proklamowanego w miejsce dawnej Rodezji Zimbabwe, 
zostawiając dwójkę małych Murzynków i ich śląską matkę w infamii.}
\poltrans{The fruits of the internationalist-sexual passions of a 
dark-skinned student, who after finishing his studies in Poland took 
himself back to the Black Continent to the newly proclaimed Zimbabwe 
in place of the former Rhodesia, leaving two small mixed-race children 
and their Silesian mother in disgrace.}

A note on language: \textbf{Murzyn} is an old Polish word for a Black
person, now considered offensive in contemporary usage. It appears here
as the narrator's reported voice, reflecting the language of the
community being described --- not as an endorsement.

\textbf{internacjonalistyczno-seksualnych uniesień} --- a compound of
the narrator's coinage, fusing two registers: \textit{internacjonalistyczny}
(internationalist) --- the official communist vocabulary of solidarity
between peoples --- and \textit{seksualny} (sexual). The hyphen joins
what the state kept carefully separate.

\textbf{czarnoskórego studenta} --- \textit{of a dark-skinned student}.
Under communism, Poland hosted many students from Africa and Asia
through solidarity agreements; this is the specific historical context.

\textbf{zabrał się z powrotem na Czarny Ląd} --- \textit{took himself
back to the Black Continent}. \textbf{Czarny Ląd} --- the traditional
Polish name for Africa.

\textbf{do proklamowanego w miejsce dawnej Rodezji Zimbabwe} --- \textit{to
the newly proclaimed Zimbabwe in place of the former Rhodesia}. The
narrator specifies the historical moment precisely: Zimbabwe was
proclaimed in 1980, dating the story and giving the father's departure
its political context --- he was returning to a newly independent nation.

\textbf{zostawiając\ldots w infamii} --- \textit{leaving\ldots in
disgrace}. \textbf{Infamia} --- public disgrace, loss of good name,
from Latin. In a small Silesian town in 1980, an unmarried woman with
two mixed-race children, abandoned, would have faced exactly this.

The sentence is long because the situation is long: a whole life
compressed into a subordinate clause. Adam and Robert grew up here,
went to school here, are buried here. Their father returned to a newly
free Zimbabwe; their mother stayed in infamy.

% ────────────────────────────────────────────────────────────

\poltext{Nie sprawdziło się czarnowidztwo śląskiego klasyka, Janoscha, 
który pisał o tych stronach: „Bo żaden Murzyn by się tu nie uchował. 
Każdego odmieńca ukradkiem nocą by zatłukli".}
\poltrans{The dark prophecy of the Silesian classic, Janosch, who wrote 
about these parts, did not come true: `Because no Black man would 
survive here. They would beat every outsider to death in secret, 
at night.'}

\textbf{Nie sprawdziło się} --- \textit{it did not come true, it was
not borne out}. \textbf{Sprawdzić się} --- to prove true. The negation
comes first and is the whole point of the sentence.

\textbf{czarnowidztwo} --- \textit{dark prophecy, doom-saying}. From
\textit{czarny} (black) + \textit{widzieć} (to see): literally
\textit{black-seeing}, the habit of seeing only the darkest possible
outcome. A rich compound.

\textbf{Janosch} --- pen name of Horst Eckert (born 1931), a writer
and illustrator born in Zabrze in what was then German Silesia, best
known internationally for his children's books but also the author of
grimly funny, unsentimental prose about Silesian life. A
\textit{klasyk} of this specific region.

\textbf{odmieńca} --- \textit{outsider, deviant, one who is different}.
From \textit{odmienny} (different, other). Stronger than mere
\textit{obcy} (stranger): an \textit{odmieniec} is someone who
deviates from the local norm, who is visibly, fundamentally other.

The quoted words are Janosch's, presented in order to be refuted ---
the graves of Adam and Robert Mitwerandu are the refutation. The
narrator doesn't argue with Janosch; he simply points at the evidence.
Two Black Silesian brothers, buried in their hometown cemetery,
remembered by their classmates with a bottle. The \textit{obchód} is
doing more than commemorating the dead; it is, in its unsteady way,
a testimony.

% ── Bracia Mitwerandu / The Mitwerandu Brothers ──────────────
\section{Bracia Mitwerandu / The Mitwerandu Brothers}

\poltext{Na przekór temu, chłopcy dobrze radzili sobie w dzieciństwie 
pośród śląskich rówieśników ze szkolnego boiska.}
\poltrans{Despite this, the boys managed well in childhood among their 
Silesian peers from the school playground.}

\textbf{Na przekór temu} --- \textit{despite this, in defiance of this}.
\textbf{Na przekór} is a fixed prepositional phrase meaning \textit{against,
contrary to, in spite of} --- with a slight flavour of active defiance
rather than passive endurance. \textbf{Temu} points back at everything
just said: the infamy, the abandonment, Janosch's prediction.

\textbf{dobrze radzili sobie} --- \textit{managed well, coped well}.
\textbf{Radzić sobie} --- \textit{to manage, to cope, to get along}
--- one of the most useful and idiomatic Polish constructions. The
reflexive \textit{sobie} is inseparable from it.

\textbf{pośród śląskich rówieśników} --- \textit{among their Silesian
peers}. \textbf{Pośród} --- \textit{among, amidst} --- slightly more
literary than \textit{wśród}. \textbf{Rówieśnicy} --- \textit{peers,
contemporaries, age-mates} --- a precise and rather beautiful word.

% ────────────────────────────────────────────────────────────

\poltext{Być może pomagała w tym siła pięści i kopniaków Roberta. 
Poza tym dwójce trudniej było spuścić łomot niż jakiejś pojedynczej 
sierotce.}
\poltrans{Perhaps Robert's fists and kicks helped with that. Besides, 
it was harder to give a beating to a pair than to some lone orphan.}

\textbf{siła pięści i kopniaków} --- \textit{the force of fists and
kicks}. \textbf{Pięść} --- \textit{fist}; \textbf{kopniak} ---
\textit{kick}. Robert was the physically formidable one; his reputation
on the playground provided a practical solution to Janosch's predicted
violence.

\textbf{spuścić łomot} --- \textit{to give a beating, to thrash}. A
vivid colloquial expression: \textbf{łomot} is the sound and act of a
heavy blow; \textbf{spuścić} (to deliver) makes the phrase literal and
physical.

\textbf{niż jakiejś pojedynczej sierotce} --- \textit{than to some
lone orphan}. \textbf{Jakiejś} --- \textit{some, any old} --- genitive
of the indefinite \textit{jakiś}, dismissive in register.
\textbf{Sierotce} --- dative of \textit{sierotka}, diminutive of
\textit{sierota} (orphan): the diminutive making it more pitiful, more
defenceless. The sentence is bleakly practical: two children together
were simply a harder target.

% ────────────────────────────────────────────────────────────

\poltext{W każdym razie trzymali się razem i w podstawówce, zamiast 
ofiarami rasizmu, częściej byli obiektami podziwu, zwłaszcza wśród 
koleżanek. Choć nie tylko młodociane adoratorki były pod ich urokiem.}
\poltrans{In any case they stuck together at primary school too, and 
instead of being victims of racism, they were more often objects of 
admiration, especially among the girls. Though it wasn't only their 
youthful admirers who were under their spell.}

\textbf{trzymali się razem} --- \textit{they stuck together, they kept
together}. A warm, solid expression for solidarity.

\textbf{zamiast ofiarami rasizmu} --- \textit{instead of being victims
of racism}. \textbf{Zamiast} + instrumental: \textit{instead of}.
\textbf{Ofiarami} --- instrumental plural of \textit{ofiara} (victim).
The word \textit{rasizm} is stated plainly and without deflection.

\textbf{obiektami podziwu} --- \textit{objects of admiration}.
\textbf{Obiektami} --- instrumental plural of \textit{obiekt}: slightly
clinical, deliberately so --- \textit{objects} rather than
\textit{subjects} of admiration, the gaze of others still constituting
them as something to be looked at, even when that look is favourable.

\textbf{adoratorki} --- \textit{female admirers, devotees} --- a
slightly old-fashioned, romantic word. \textbf{Pod urokiem} ---
\textit{under the spell/charm of} --- a fixed expression.
\textbf{Urok}: charm, spell, enchantment --- the same word used in
fairy tales and in everyday compliments. The final sentence opens a
door that the next passage walks through.

% ────────────────────────────────────────────────────────────

\poltext{--- Zobocz, jak się oni gryfnie nauczyli godać! --- pochwaliła. 
Dopiero gdy z ust chłopaków zaczęły lecieć kurwy i chuje, dodała z 
ledwo wyczuwalną już nutką podziwu: --- A chachory z nich som\ldots}
\poltrans{`See how nicely they've learned to talk!' --- she said 
approvingly. Only when expletives started flying from the boys' mouths 
did she add, with a barely perceptible note of admiration: `And they're 
proper rascals, they are\ldots'}

This passage is dense with Silesian dialect --- arguably the richest
single sentence in the chapter.

\textbf{Zobocz} --- Silesian dialect for \textit{patrz} or
\textit{spójrz} (look, see). The specifically Silesian imperative form,
immediately placing the speaker within the regional vernacular.

\textbf{gryfnie} --- \textit{nicely, beautifully, well}. One of the
most characteristic Silesian dialect words, from the German
\textit{griffig} (handy, neat). In standard Polish: \textit{ładnie} or
\textit{pięknie}. \textit{Gryfnie} is warm and approving, with a
slightly old-fashioned tenderness.

\textbf{godać} --- \textit{to speak, to talk}. The quintessential
Silesian dialect verb, equivalent to standard Polish \textit{mówić}.
To \textit{godać} in Silesia is not merely to speak but to speak
\textit{Silesian} --- it is the word the dialect uses to describe
itself. The boys have been admitted into the linguistic community of
the region.

\textbf{z ust chłopaków zaczęły lecieć kurwy i chuje} --- expletives
\textit{started flying} from the boys' mouths. \textbf{Lecieć} (to
fly) is the standard colloquial verb for language that erupts freely.
The two expletives are the precise linguistic achievement that earned
the boys their final accolade --- fluency in a language includes its
profanity.

\textbf{z ledwo wyczuwalną już nutką podziwu} --- \textit{with a barely
perceptible note of admiration}. \textbf{Ledwo wyczuwalną} ---
\textit{barely detectable}. \textbf{Nutka} --- \textit{a little note,
a hint} --- diminutive of \textit{nuta} (musical note). The admiration
is present but restrained, as if she doesn't quite want to admit it.

\textbf{A chachory z nich som} --- \textit{And they're proper rascals,
they are}. Every word Silesian dialect. \textbf{Chachory} --- rascals,
scamps, rogues: not condemnation but rough, affectionate approval for
someone spirited and irreverent. \textbf{Som} --- Silesian dialect for
\textit{są} (they are). She began praising them in Silesian and ends
her verdict in the same tongue.

The ultimate credential for belonging was not merely speaking the
dialect --- that earned \textit{gryfnie} --- but swearing fluently in
it. The \textit{kurwy i chuje} are the final examination, and the boys
passed. \textit{Chachory z nich som} is an induction: you are one of
ours now.

% ── Robert / Robert ──────────────────────────────────────────
\section{Robert}

\poltext{Groszek zaczął wspominać: --- Robert nawet dobrze się 
zapowiadał jako zawodnik. Miał talent. Grał chyba w okręgówce, ale 
miał przenieść się do pierwszej ligi. Nie pamiętam --- do Gieksy?}
\poltrans{Groszek began to reminisce: `Robert even showed real promise 
as a player. He had talent. He played in the regional league I think, 
but he was due to move up to the first division. I don't remember --- 
to GKS?'}

\textbf{zaczął wspominać} --- \textit{began to reminisce}.
\textbf{Wspominać} --- imperfective: to reminisce, to recall. The
imperfective signals an open-ended flow of memories beginning. They
are standing at Robert's grave; the marble ball now has context.

\textbf{dobrze się zapowiadał jako zawodnik} --- \textit{showed real
promise as a player}. \textbf{Zapowiadać się} --- \textit{to show
promise, to look promising} --- literally \textit{to announce oneself},
to signal what one might become.

\textbf{Grał chyba w okręgówce} --- \textit{he played in the regional
league I think}. \textbf{Chyba} --- \textit{I think, probably} ---
the particle of uncertain recollection, very common in spoken Polish.
\textbf{Okręgówka} --- the colloquial name for the \textit{liga
okręgowa}, the regional/district football league: several tiers below
the top flight.

\textbf{miał przenieść się do pierwszej ligi} --- \textit{he was due
to move up to the first division}. \textbf{Mieć} + infinitive: he was
supposed to, the move was anticipated or arranged.

\textbf{Gieksa} --- the affectionate colloquial name for GKS Katowice
(\textit{Górniczy Klub Sportowy} --- the Miners' Sports Club), one of
the major Silesian football clubs. The name is formed from the letters
G-K-S sounded out, then given the feminine diminutive ending: a club
name made tender by familiarity.

The marble ball on the grave suddenly resolves. It wasn't a quirky
choice --- it was a football. Robert Mitwerandu was a footballer, and
his grave is marked accordingly.

% ────────────────────────────────────────────────────────────

\poltext{--- Nie tylko, grał w ekstraklasie siedem razy, a nawet 
powołano go do reprezentacji Polski juniorów --- Waldi był w kwestiach 
futbolu dobrze zorientowany.}
\poltrans{`Not only that --- he played in the top flight seven times, 
and was even called up to the Polish under-21 squad' --- Waldi was 
well informed on football matters.}

\textbf{grał w ekstraklasie siedem razy} --- \textit{he played in the
Ekstraklasa seven times}. \textbf{Ekstraklasa} --- the top tier of
Polish professional football. Not \textit{pierwsza liga} (which
Groszek hazarded) but \textit{ekstraklasa} --- the highest level.
\textbf{Siedem razy} --- \textit{seven times}: not a long career, but
real appearances, counted and remembered.

\textbf{powołano go do reprezentacji Polski juniorów} --- \textit{he
was called up to the Polish under-21 squad}. \textbf{Powołano} ---
impersonal passive perfective: the official verb for selection to a
national squad. Robert Mitwerandu, Black Silesian, \textit{chachar},
marble-ball-on-the-grave, was called up to represent Poland.

\textbf{Waldi był w kwestiach futbolu dobrze zorientowany} ---
\textit{Waldi was well informed on football matters}. \textbf{W
kwestiach} --- \textit{in matters of} --- a slightly formal
construction giving Waldi's knowledge quasi-expert status. Waldi, on
crutches, \textit{serwis pogwarancyjny}, supplies the precise details
that Groszek's affectionate memory had blurred.

The correction matters enormously. Robert was not merely promising
--- he achieved. The marble ball is not a symbol of what might have
been but of what was.

% ────────────────────────────────────────────────────────────

\poltext{--- I nagle umarł. Podobno na serce. Wada wrodzona? Zawał? 
Doping? Nie wiadomo --- dodała moja siostra, krzywiąc się od 
cmentarnej wyborowej.}
\poltrans{`And then he suddenly died. Apparently of the heart. A 
congenital defect? A heart attack? Doping? Nobody knows' --- my sister 
added, wincing from the cemetery vodka.}

\textbf{I nagle umarł} --- the same \textbf{nagle} from the chapter's
opening sentence (\textit{Fiszer umarł nagle, niespodziewanie i
kłopotliwie}). Here it returns starker, with no companions.

\textbf{Podobno na serce} --- \textit{apparently of the heart}.
\textbf{Podobno} --- \textit{apparently, reportedly, so they say} ---
the particle of unverified information. \textbf{Na serce}: the
preposition \textit{na} with the accusative is used for causes of
death in Polish (\textit{umrzeć na raka} --- to die of cancer).

\textbf{Wada wrodzona? Zawał? Doping?} --- three questions, three
bare nouns. \textbf{Zawał} --- \textit{heart attack}, short for
\textit{zawał serca}. \textbf{Doping?} --- the question that changes
the register: the possibility that the means used to reach the
ekstraklasa may have contributed to what ended him there. No answer
offered.

\textbf{Nie wiadomo} --- \textit{nobody knows, it is not known}. The
impersonal closing of all three questions at once.

\textbf{krzywiąc się od cmentarnej wyborowej} --- \textit{wincing
from the cemetery vodka}. \textbf{Krzywić się} --- to wince, to
grimace. \textbf{Wyborowa} is a well-known Polish vodka brand (the
name means \textit{choice, select, premium}); \textbf{cmentarna} is
the narrator's adjective --- the bottle Groszek produced at the grave,
drunk \textit{z gwinta} with tears in eyes, is still going. The
narrator's sister, who first appeared explaining Sławek's absence
(\textit{wyjaśnia siostra}), delivers the final unknowable fact about
Robert while wincing from it.

% ── Adam / Adam ──────────────────────────────────────────────
\section{Adam}

\poltext{--- A Adam? Dlaczego Adam? --- zapytałem. Groszek miał z nim 
kontakt na studiach: --- On się nigdy nie umiał odnaleźć. Był wybitnie 
zdolny na polibudzie. Z matematyki. Ale to był raczej typ filozofa. 
Kiedy miał dobre dni, dziewczyny za nim szalały. Kiedy miał złe, 
lądował na Korczaka, u czubków.}
\poltrans{`And Adam? Why Adam?' --- I asked. Groszek had been in 
contact with him at university: `He never managed to find himself. He 
was outstandingly gifted at the polytechnic. In mathematics. But he 
was more of a philosopher type. When he had good days, the girls were 
crazy about him. When he had bad ones, he ended up on Korczaka, with 
the loonies.'}

\textbf{On się nigdy nie umiał odnaleźć} --- \textit{he never managed
to find himself}. \textbf{Umiał} --- past of \textit{umieć} (to know
how to, to be capable of) --- not \textit{mógł} (was able to, had the
opportunity): he lacked the capacity itself. \textbf{Odnaleźć się}
--- \textit{to find oneself, to locate oneself in life} --- not merely
\textit{to find one's place} but to successfully inhabit one's own
existence.

\textbf{wybitnie zdolny na polibudzie} --- \textit{outstandingly gifted
at the polytechnic}. \textbf{Wybitnie} is a strong intensifier, beyond
\textit{bardzo} (very). \textbf{Na polibudzie} --- locative of
\textit{polibuda}, the affectionate colloquial name for the
\textit{politechnika} --- the same institution that produced Groszek,
who delivers this judgement with the authority of a fellow alumnus.

\textbf{raczej typ filozofa} --- \textit{more of a philosopher type}.
\textbf{Raczej} --- \textit{rather, more of a}. The mathematical gift
was real, but it sat inside someone whose nature pulled toward
abstraction and inwardness.

\textbf{dziewczyny za nim szalały} --- \textit{the girls were crazy
about him}. \textbf{Szaleć za kimś} --- \textit{to be mad for someone}
--- \textit{szaleć} is \textit{to rage, to go mad}: the girls'
attraction described as productive madness.

\textbf{lądował na Korczaka, u czubków} --- \textit{he ended up on
Korczaka, with the loonies}. \textbf{Lądował} --- \textit{he would
land} --- imperfective, the repeated landing: this happened more than
once. \textbf{Na Korczaka} --- \textit{on Korczaka Street}: the
psychiatric hospital known locally by its address. \textbf{U czubków}
--- \textit{with the loonies}: \textbf{czubek} (colloquially, a mad
person) from \textit{czub} (a tuft, a crest). Groszek's word, spoken
with rough affection.

Adam's portrait is heartbreaking in its symmetry with Robert's. Robert
had the body --- the ekstraklasa, the marble ball, the heart that
failed. Adam had the mind --- the mathematics, the philosophy, the
girls, the psychiatric ward. \textit{On się nigdy nie umiał odnaleźć}:
he had everything except the ability to be at home in his own life.

% ────────────────────────────────────────────────────────────

\poltext{(Dopisek topograficzno-historyczny: To adres lokalnych Tworek, 
w bliskim sąsiedztwie wieży wodnej. Stary budynek szpitala, tak jak 
wieżę, zaprojektowali Zillmannowie. Teraz zrujnowany, w czasach 
świetności stylem przypominał sopocki Grand Hotel, tyle że zamiast na 
molo wychodziło się z niego na końskie doły, dawne wyrobiska po piasku 
do kopalni. O ile się wychodziło, bo z psychiatryka na ogół łatwo nie 
wypuszczali).}
\poltrans{(A topographic-historical note: This is the address of the 
local Tworki, in the close vicinity of the water tower. The old hospital 
building, like the tower, was designed by the Zillmanns. Now ruined, in 
its heyday it was reminiscent in style of the Grand Hotel in Sopot, 
except that instead of onto a pier you stepped out onto horse pits, 
former sand extraction workings for the mine. If you stepped out at all, 
since they didn't generally let you out of a psychiatric hospital 
easily.)}

\textbf{Dopisek topograficzno-historyczny} --- \textit{a
topographic-historical note}. \textbf{Dopisek} --- literally \textit{a
written addition}, a note appended. The compound adjective is the
narrator playing the amateur local historian --- the same impulse that
produced \textit{cmentarny opis obyczajów minionego wieku}.

\textbf{Tworki} --- a colloquial name for a famous psychiatric hospital
outside Warsaw, used here as a common noun: \textit{the local Tworki}
--- the local equivalent. Every Polish town applies this name to its
own facility by analogy with the nationally infamous original.

\textbf{zaprojektowali Zillmannowie} --- \textit{was designed by the
Zillmanns}. A dynasty of architects, German in name, active in Silesia,
responsible for both the hospital and the water tower. Named with the
casual confidence of someone who knows local architectural history.

\textbf{w czasach świetności stylem przypominał sopocki Grand Hotel}
--- \textit{in its heyday reminiscent in style of the Grand Hotel in
Sopot}. The Grand Hotel in Sopot is a landmark of Polish seaside
elegance, a belle époque pile on the Baltic. The comparison is real:
both late nineteenth-century, both grand in ambition.

\textbf{zamiast na molo wychodziło się z niego na końskie doły} ---
\textit{instead of onto a pier you stepped out onto horse pits}.
\textbf{Molo} --- \textit{pier, jetty} --- Sopot's defining feature.
\textbf{Końskie doły} --- \textit{horse pits}: excavation pits named
for the horses that hauled extracted material. The same architectural
grandeur, a completely different view from the terrace.

\textbf{dawne wyrobiska po piasku do kopalni} --- \textit{former sand
extraction workings for the mine}. Sand mined industrially for the
coal mine. Silesian landscape in miniature.

\textbf{O ile się wychodziło} --- \textit{if you stepped out at all}
--- the parenthesis closes on this. The grandeur of the view was, in
any case, largely theoretical for the hospital's residents. The
architecture of aspiration and the landscape of extraction --- the same
Silesian juxtaposition the cemetery survey observed --- and in the
middle of it, Adam, brilliant and lost.


% Subsequent chapters — uncomment as work progresses
% % ============================================================
%  ch02_ostatni_dzien_wakacji.tex  —  Chapter 2
%  Ostatni dzień wakacji. Osiedle na Burowcu (pp. TBC)
% ============================================================
%
% ── STYLE RULES FOR THIS FILE ───────────────────────────────
%
%  These rules must be followed consistently throughout.
%  Drift from these rules has occurred before — check carefully.
%
%  HEADINGS:
%  - Use \section{} for ALL passage headings without exception.
%  - \subsection{} and \subsection*{} are NOT used in this file.
%  - Section titles follow the pattern:
%      Polish title / English title   (e.g. \section{Obchód / The Rounds})
%    or just a name if self-evident   (e.g. \section{Groszek})
%  - \section{} renders in blue bold with a rule beneath
%    (defined in olowiane_dzieci_analiza.tex).
%
%  POLISH SOURCE TEXT:
%  - Always use \poltext{...} — never \begin{quote}...\end{quote}
%    and never plain \textit{...} or any other wrapper.
%  - \poltext{} renders the text in red italic (keywordred).
%
%  ENGLISH TRANSLATION:
%  - Always use \poltrans{...} immediately after \poltext{}.
%  - \poltrans{} renders in blue (polishblue) with a rule beneath.
%  - No label precedes the translation — colour contrast suffices.
%
%  ANALYSIS NOTES:
%  - Follow \poltrans{} with free prose commentary.
%  - Bold Polish terms with \textbf{}, gloss in \textit{}.
%  - Standard pattern: \textbf{term} --- \textit{gloss} --- analysis.
%  - Use \medskip before thematic summary paragraphs.
%  - Grammar or vocabulary notes may be prefixed with:
%    \noindent\textit{Grammar note:} or \noindent\textit{Vocabulary note:}
%
%  PASSAGE SEPARATORS:
%  - Use % ──────────────────────────────────────────────────
%    (single % with em-dashes) between passages within a section.
%    Do NOT use the %% --- style.
%
%  UPDATE PROTOCOL:
%  - The LaTeX file is updated after each passage is confirmed
%    by the user — not in batches.
%  - Always append new content by str_replace on the final line.
%  - Verify line count after each update with wc -l.
%
% ── WHERE THE READING CURRENTLY STANDS ─────────────────────
%
%  Last passage analysed (session ending 2026-03-06):
%  Chapter 2 not yet started.
%  Chapter 1 (ch01_jeszcze_jeden_pogrzeb.tex) is complete.
%
% ============================================================


\chapter{Ostatni dzień wakacji. Osiedle na Burowcu}
\markboth{Ostatni dzień wakacji. Osiedle na Burowcu}{}

% ── Chapter title ────────────────────────────────────────────
\section{Tytuł rozdziału / Chapter Title}

\poltext{Ostatni dzień wakacji. Osiedle na Burowcu.}
\poltrans{The Last Day of the Holidays. The Estate at Burowiec.}

\textbf{Ostatni dzień wakacji} --- \textit{the last day of
the holidays}. \textbf{Ostatni} --- \textit{last, final} ---
from \textit{ostać się} (to remain); the superlative of
lateness, the day at the end of summer that everyone
who has been a schoolchild knows by feel.
\textbf{Wakacje} --- \textit{school holidays, summer
vacation} --- plural only, like \textit{drzwi} (door)
or \textit{nożyczki} (scissors): summer does not come
in the singular. The word is borrowed from Latin
\textit{vacatio} (freedom, exemption) via the academic
calendar.

\textbf{Osiedle na Burowcu} --- \textit{the estate at
Burowiec}. \textbf{Osiedle} --- the housing estate, the
\textit{bloki}: the word that has governed the spatial
world of Chapter~1. \textbf{Na Burowcu} --- locative
of \textbf{Burowiec}: a specific named estate or district,
\textit{na} indicating location on or at a place
(used with open areas, estates, squares, as distinct
from \textit{w} for enclosed spaces). Burowiec is a
district of Olkusz, a town in the Małopolska region
of southern Poland, southeast of Kraków --- the area
whose proximity to lead and zinc smelting operations
is the book's subject.

\medskip
Where Chapter~1 opened at a funeral and moved through
the present-tense reckoning of a generation's losses,
Chapter~2 announces its coordinates in time and space:
the last day of summer, a specific estate in a specific
town. The chapter will move into the past, into the
children themselves --- before the vanishing, at the
moment just before school resumes and the world of the
estate is still fully inhabited.

%% ---------------------------------------------------------------
\section{Wszystko zaczęło się od huty / It All Started with the Smelter}
%% ---------------------------------------------------------------

\poltext{Co do znikania moich rówieśników, wszystko zaczęło
się od huty. Przez sto pięćdziesiąt lat jej istnienia różnie
ją nazywano: Uthemann, Zakłady Cynkowe Szopienice, blajówa,
Huta Metali Nieżelaznych. Jak zwał, tak zwał. Tu wszystko
zaczynało i kończyło się na hucie.}

\poltrans{As for the disappearing of my contemporaries,
it all started with the smelter. Over a hundred and fifty
years of its existence it was called by various names:
Uthemann, the Szopienice Zinc Works, \textit{blajówa},
the Non-Ferrous Metals Smelter. Whatever it was called,
it was called. Here everything began and ended with the
smelter.}

The chapter opens with a declaration, a list, a proverb,
and an axiom. In four sentences Jędryk names the smelter
as the origin of everything, passes through 150 years of
its history in a list of aliases, dismisses the question
of names with a shrug, and then restates the first claim
with the finality of a closing argument. The book has
found its subject; Chapter~2 announces it without ceremony.

\textbf{Co do znikania moich rówieśników} ---
\textit{as for the disappearing of my contemporaries}.
\textbf{Co do} + genitive: a formal topic-introducer,
\textit{as for, regarding, on the matter of} ---
the register of a deposition or a report. \textbf{Znikania}
--- genitive of \textit{znikanie}, the verbal noun from
\textit{znikać}: \textit{the disappearing}, the process
nominalised. Chapter~1 used \textit{znikały} and
\textit{znikali} as verbs; now the disappearing is a
noun, a phenomenon to be accounted for. \textbf{Moich
rówieśników} --- \textit{of my contemporaries, of my
peers}: genitive plural of \textit{rówieśnik} (from
\textit{równy} --- equal, and \textit{wiek} --- age):
one who is of equal age, a peer. \textbf{Moich} ---
\textit{my} --- the possessive confirms the narrator's
place inside the cohort: not an external observer but
a member of the disappearing generation.

\textbf{wszystko zaczęło się od huty} ---
\textit{it all started with the smelter}. The thesis.
\textbf{Wszystko} --- \textit{everything, it all}.
\textbf{Zaczęło się} --- perfective reflexive past of
\textit{zaczynać się}: \textit{it began, it started}
--- the reflexive \textit{się} makes the beginning
intransitive, self-initiated, as if the thing started
of its own accord. \textbf{Od huty} --- \textit{from
the smelter}: \textit{od} + genitive of origin and
cause. The smelter is the source, the point of
departure, the \textit{od} from which all consequences
flow. The sentence could close the book as easily as
open a chapter.

\textbf{Przez sto pięćdziesiąt lat jej istnienia} ---
\textit{over a hundred and fifty years of its existence}.
\textbf{Przez} + accusative: duration, passage through
time. \textbf{Sto pięćdziesiąt lat} --- one hundred
and fifty years. \textbf{Jej istnienia} --- genitive
of \textit{jej istnienie}: \textit{its existence},
the smelter's span of operation. A century and a half:
the smelter predates the narrator's grandparents,
predates modern Poland, predates two world wars.

\textbf{różnie ją nazywano} --- \textit{it was called
by various names, people called it variously}.
\textbf{Różnie} --- \textit{variously, in different
ways}. \textbf{Ją} --- accusative of \textit{ona}:
\textit{her/it}, referring to \textit{huta} (feminine).
\textbf{Nazywano} --- impersonal passive past of
\textit{nazywać}: \textit{one called, it was called}
--- the unnamed collective of namers, the succession
of owners, regimes, and bureaucracies that renamed
it over 150 years.

\medskip
The list of names that follows is a compressed history
of the region:

\textbf{Uthemann} --- the original founder's name:
Friedrich Uthemann, a German industrialist who
established a zinc-smelting operation in Szopienice
(now a district of Katowice, in Upper Silesia) in
the nineteenth century. The German surname places the
smelter's origins in the industrial colonisation of
Silesia under Prussian rule.

\textbf{Zakłady Cynkowe Szopienice} --- \textit{the
Szopienice Zinc Works}: the interwar or early
communist-era designation, a sober administrative
title. \textbf{Zakłady} --- \textit{works, plant,
establishment} (plural of \textit{zakład}).
\textbf{Cynkowe} --- \textit{zinc} (adjectival form
of \textit{cynk}).

\textbf{blajówa} --- the workers' name, the local
colloquial name, set deliberately among the official
designations. From German \textit{Blei} (lead) +
the Polish suffix \textit{-ówa} (used to form
informal place-names and nicknames, especially
feminine): \textit{the lead-place, the leady one}.
The suffix makes it familiar, almost affectionate ---
the name the neighbourhood used, not the name on
the letterhead. Its presence in this list is the
list's most important entry: it is the only name
that tells you what was actually produced there.

\textbf{Huta Metali Nieżelaznych} --- \textit{the
Non-Ferrous Metals Smelter}: the late communist-era
official designation, bureaucratically precise and
informationally evasive. \textbf{Nieżelaznych} ---
\textit{non-ferrous}: genitive plural of
\textit{nieżelazny} (\textit{nie-} + \textit{żelazny},
iron): not iron, i.e.\ lead, zinc, cadmium. The
name describes the product category without naming
any specific metal. The more official the name, the
less it says.

\textbf{Jak zwał, tak zwał.} --- \textit{Whatever
it was called, it was called. A name's a name.}
A Polish proverb-formula: \textbf{jak\ldots{}
tak\ldots{}} sets up a parallelism of indifference.
\textbf{Zwał} is the past tense third person
singular of \textit{zwać} --- an archaic/literary
verb meaning \textit{to call, to name} --- used
here as the impersonal \textit{one called it}.
The formula means: \textit{it doesn't matter what
you called it; the thing was what it was}. The
proverb dismisses the whole question of naming
that the preceding list raised, and in doing so
makes a point: the names changed (German founder,
zinc works, workers' slang, communist designation),
but the smelter and its emissions did not.

\textbf{Tu wszystko zaczynało i kończyło się na
hucie.} --- \textit{Here everything began and ended
with the smelter.} The chapter's axiom, restating
the opening thesis with added weight. \textbf{Tu}
--- \textit{here}: the estate, the neighbourhood,
the world of these children. \textbf{Zaczynało i
kończyło się} --- imperfective reflexive past:
\textit{began and ended}, the two verbs paired and
balanced, \textit{zaczynać się} (to begin) and
\textit{kończyć się} (to end). \textbf{Na hucie}
--- \textit{at/with the smelter}: locative of
\textit{huta}, the \textit{na} of attachment and
dependence. Everything in this place had its origin
and its terminus in the smelter: employment, economy,
community --- and, the book argues, illness and death.

\medskip
\noindent\textit{Context note:} The smelter in
question --- operating under its various names in
Szopienice and the surrounding Silesian--Małopolska
region --- processed lead and zinc ore for over a
century, emitting lead particulates that settled on
the surrounding estates. The children who grew up
there in the 1970s and 1980s were breathing and
ingesting lead dust at levels now understood to
cause neurological damage, immune dysfunction, and
early death. The \textit{blajówa} --- the lead-place
--- named itself honestly. The official designations
did not.

\medskip
The opening of Chapter~2 performs the same movement
as the close of Chapter~1, but inverted: where
\textit{znikającą klasę} arrived at the end as a
discovery, \textit{wszystko zaczęło się od huty}
is stated at the start as a premise. The chapter
will now go back to show how the premise was lived.

% \input{ch03_krolowa}
% \input{ch04_znikajaca_klasa}
% \input{ch05_diagnoza_szwajcaria}
% \input{ch06_pacjent_zero}
% \input{ch07_osiedlowy_detektyw}
% \input{ch08_rob_co_chcesz}
% \input{ch09_teenage_dream}
% \input{ch10_krolowa_wojewoda_i_cysorz}
% \input{ch11_placz_pod_wieza}
% \input{ch12_boj_w_hucie}
% \input{ch13_szklana_niepogoda}
% \input{ch14_exodus_od_smrodu}
% \input{ch15_badanie_glowy}
% \input{ch16_szpieg_i_ortografia}
% \input{ch17_dwa_razy_kwestia_mieszkaniowa}
% \input{ch18_pozegnanie}
% \input{ch19_trzeba_isc_z_tym_wyzej}
% \input{ch20_dzien_w_podrozy}
% \input{ch21_radosc_na_placu_dzierzynskiego}
% \input{ch22_na_przelaczy}
% \input{ch23_strefa_ochronna}
% \input{ch24_babskie_wembley}
% \input{ch25_daleko_od_domu}
% \input{ch26_matka_w_zielonym_fiacie}
% \input{ch27_honoris_causa}
% \input{ch28_powroty}
% \input{ch29_gorzka_prawda}
% \input{ch30_dwie_matki}
% \input{ch31_izba_tradycji}
% \input{ch32_olowiany_balast}
% \input{ch33_od_autora}

\end{document}
